%==============================================================================
%  APOSTILA DE CONTROLE — BALL BALANCER 3RPS (ESP32-S3)
%  Material didatico completo: da modelagem fisica ao controle digital avancado.
%  Compila em Overleaf com pdfLaTeX. Passes: pdflatex -> bibtex -> pdflatex x2
%  -> makeindex -> pdflatex.  (No Overleaf, basta apertar "Recompile".)
%==============================================================================

% --- Base de referencias embutida (BibTeX em arquivo unico) -------------------
\begin{filecontents}[overwrite]{referencias.bib}
@book{ogata,
  author    = {Ogata, Katsuhiko},
  title     = {Engenharia de Controle Moderno},
  edition   = {5},
  publisher = {Pearson Prentice Hall},
  year      = {2010}
}
@book{nise,
  author    = {Nise, Norman S.},
  title     = {Engenharia de Sistemas de Controle},
  edition   = {7},
  publisher = {LTC},
  year      = {2017}
}
@book{franklin,
  author    = {Franklin, Gene F. and Powell, J. David and Workman, Michael},
  title     = {Digital Control of Dynamic Systems},
  edition   = {3},
  publisher = {Addison-Wesley},
  year      = {1998}
}
@book{astrom,
  author    = {{\AA}str{\"o}m, Karl J. and Murray, Richard M.},
  title     = {Feedback Systems: An Introduction for Scientists and Engineers},
  publisher = {Princeton University Press},
  year      = {2008}
}
@phdthesis{cordeiro,
  author = {Cordeiro, Giovanna Ferreira},
  title  = {Controle Cl{\'a}ssico da Planta Ball Balancer},
  school = {UNESP/FEIS, Ilha Solteira},
  year   = {2022},
  type   = {Trabalho de Gradua{\c c}{\~a}o em Engenharia El{\'e}trica}
}
@misc{musa,
  author       = {Musa, Aaed},
  title        = {Ball Balancer},
  howpublished = {Instructables},
  year         = {2021},
  note         = {Dispon{\'i}vel em: https://www.instructables.com/Ball-Balancer/}
}
@manual{esp32s3,
  author       = {{Espressif Systems}},
  title        = {ESP32-S3 Technical Reference Manual},
  year         = {2023}
}
\end{filecontents}

\documentclass[11pt,oneside,a4paper]{book}

% --- Codificacao e idioma -----------------------------------------------------
\usepackage[utf8]{inputenc}
\usepackage[T1]{fontenc}
\usepackage[brazil]{babel}
\usepackage{lmodern}
\usepackage{microtype}

% --- Matematica ---------------------------------------------------------------
\usepackage{amsmath,amssymb,amsthm}
\usepackage{mathtools}

% --- Layout -------------------------------------------------------------------
\usepackage[a4paper,margin=2.6cm,headheight=14pt]{geometry}
\usepackage{xcolor}
\usepackage{graphicx}
\usepackage{booktabs}
\usepackage{tabularx}
\usepackage{array}
\usepackage{ragged2e}
\usepackage{multirow}
\usepackage{longtable}
\usepackage{enumitem}
\usepackage{caption}
\usepackage{fancyhdr}
\usepackage{titlesec}
\usepackage[most]{tcolorbox}
\usepackage{listings}
\usepackage{makeidx}

% --- Graficos e diagramas -----------------------------------------------------
\usepackage{tikz}
\usetikzlibrary{arrows.meta,positioning,calc,shapes.geometric,shapes.misc,
  decorations.pathmorphing,decorations.pathreplacing,backgrounds,fit,patterns,angles,quotes}
\usepackage{pgfplots}
\pgfplotsset{compat=1.18}

% --- Hyperref (carregar quase por ultimo) -------------------------------------
\usepackage{hyperref}

\makeindex

%==============================================================================
%  PALETA E ESTILOS
%==============================================================================
\definecolor{corprim}{HTML}{12355B}   % azul-petroleo escuro
\definecolor{corsec}{HTML}{1B998B}    % verde-agua
\definecolor{coracc}{HTML}{E07A5F}    % laranja terroso
\definecolor{coralert}{HTML}{C1121F}  % vermelho
\definecolor{corroxo}{HTML}{6A4C93}   % roxo
\definecolor{cortinta}{HTML}{1A1A2E}
\definecolor{corcinza}{HTML}{6C757D}

\hypersetup{
  colorlinks=true, linkcolor=corprim, citecolor=corsec, urlcolor=corsec,
  pdftitle={Apostila de Controle - Ball Balancer 3RPS},
  pdfauthor={Projeto Ball Balancer 3RPS}
}

% --- Cabecalho/rodape ---------------------------------------------------------
\pagestyle{fancy}
\fancyhf{}
\renewcommand{\headrulewidth}{0.4pt}
\fancyhead[L]{\small\color{corcinza}\leftmark}
\fancyhead[R]{\small\color{corcinza}Ball Balancer 3RPS}
\fancyfoot[C]{\small\color{corcinza}\thepage}
\renewcommand{\chaptermark}[1]{\markboth{\thechapter.\ #1}{}}

% --- Titulos de capitulo ------------------------------------------------------
\titleformat{\chapter}[display]
  {\normalfont\sffamily\bfseries\color{corprim}}
  {\filright\Large\color{corsec}\MakeUppercase{\chaptertitlename}\ \Huge\thechapter}
  {8pt}{\Huge\filright}
\titlespacing*{\chapter}{0pt}{-10pt}{24pt}
\titleformat{\section}{\normalfont\large\sffamily\bfseries\color{corprim}}{\thesection}{0.6em}{}
\titleformat{\subsection}{\normalfont\sffamily\bfseries\color{corsec!80!black}}{\thesubsection}{0.6em}{}

% --- Caixas didaticas (tcolorbox) ---------------------------------------------
\newtcolorbox{conceito}[1][Conceito]{enhanced,breakable,colback=corprim!4,
  colframe=corprim,coltitle=white,fonttitle=\bfseries\sffamily,title=#1,
  boxrule=0.6pt,arc=2pt,left=6pt,right=6pt,top=4pt,bottom=4pt}
\newtcolorbox{projeto}[1][No projeto Ball Balancer]{enhanced,breakable,
  colback=corsec!6,colframe=corsec!80!black,coltitle=white,
  fonttitle=\bfseries\sffamily,title=#1,boxrule=0.6pt,arc=2pt}
\newtcolorbox{exemplo}[1][Exemplo numerico]{enhanced,breakable,colback=coracc!8,
  colframe=coracc!85!black,coltitle=white,fonttitle=\bfseries\sffamily,title=#1,
  boxrule=0.6pt,arc=2pt}
\newtcolorbox{intuicao}[1][Intuicao]{enhanced,breakable,colback=corroxo!6,
  colframe=corroxo,coltitle=white,fonttitle=\bfseries\sffamily,title=#1,
  boxrule=0.6pt,arc=2pt}
\newtcolorbox{atencao}[1][Atencao]{enhanced,breakable,colback=coralert!6,
  colframe=coralert,coltitle=white,fonttitle=\bfseries\sffamily,title=#1,
  boxrule=0.6pt,arc=2pt}

% --- Codigo C -----------------------------------------------------------------
\lstdefinestyle{cstyle}{
  language=C, basicstyle=\ttfamily\footnotesize,
  keywordstyle=\color{corprim}\bfseries,
  commentstyle=\color{corsec!70!black}\itshape,
  stringstyle=\color{coracc!85!black},
  numbers=left, numberstyle=\tiny\color{corcinza}, numbersep=8pt,
  frame=single, rulecolor=\color{corcinza!50}, framerule=0.4pt,
  backgroundcolor=\color{corprim!3},
  breaklines=true, showstringspaces=false, tabsize=2, captionpos=b,
  morekeywords={float,bool,int32_t,uint32_t,static,inline}
}
\lstset{style=cstyle}

% --- Estilos TikZ para diagramas de blocos ------------------------------------
\tikzset{
  bloco/.style={draw=corprim,fill=corprim!8,rounded corners=2pt,thick,
    minimum width=2.2cm,minimum height=1cm,align=center,font=\small},
  blocoS/.style={draw=corsec!80!black,fill=corsec!12,rounded corners=2pt,thick,
    minimum width=2.2cm,minimum height=1cm,align=center,font=\small},
  soma/.style={draw=corprim,thick,circle,minimum size=7mm,inner sep=0pt},
  seta/.style={-{Latex[length=2.6mm,width=2mm]},thick},
  rotulo/.style={font=\footnotesize\itshape,color=corcinza}
}

% --- Ambientes de teorema/definicao -------------------------------------------
\theoremstyle{definition}
\newtheorem{definicao}{Defini\c{c}\~ao}[chapter]
\newtheorem{exemplon}{Exemplo}[chapter]

\newcommand{\Rs}{R\$}

%==============================================================================
\begin{document}
\frontmatter

%------------------------------------------------------------------ CAPA -------
\begin{titlepage}
\begin{tikzpicture}[remember picture,overlay]
  \fill[corprim] (current page.north west) rectangle ([yshift=-6.2cm]current page.north east);
  \fill[corsec]  ([yshift=-6.2cm]current page.north west) rectangle ([yshift=-6.7cm]current page.north east);
  % bola e plano estilizados
  \begin{scope}[shift={([yshift=-3.1cm,xshift=-3.2cm]current page.north east)}]
    \draw[white,very thick,rotate=-12] (-1.6,0) -- (1.6,0);
    \shade[ball color=coracc,rotate=-12] (-0.6,0.35) circle (0.45);
  \end{scope}
  \node[anchor=north west,text=white,font=\sffamily\bfseries\Huge]
    at ([yshift=-1.4cm,xshift=1.6cm]current page.north west) {APOSTILA DE CONTROLE};
  \node[anchor=north west,text=white,font=\sffamily\large]
    at ([yshift=-2.7cm,xshift=1.6cm]current page.north west) {Do duplo integrador ao controle digital};
  \node[anchor=north west,text=corsec!20!white,font=\sffamily\large]
    at ([yshift=-3.5cm,xshift=1.6cm]current page.north west) {aplicado ao Ball Balancer 3RPS (ESP32-S3)};

  \fill[corprim] (current page.south west) rectangle ([yshift=2.2cm]current page.south east);
  \node[anchor=south,text=white,font=\sffamily\small]
    at ([yshift=1.0cm]current page.south) {Engenharia de Controle $\cdot$ Controle Digital $\cdot$ Sistemas Embarcados};
  \node[anchor=south,text=corsec!30!white,font=\sffamily\footnotesize]
    at ([yshift=0.5cm]current page.south) {Material did\'atico do projeto $\cdot$ \today};
\end{tikzpicture}

\vspace*{7.5cm}
\noindent\begin{tcolorbox}[colback=white,colframe=corsec,boxrule=1pt,arc=3pt]
\sffamily\large
Uma apostila que parte das \textbf{leis de Newton} e chega ao \textbf{PID discreto
embarcado}, sempre conectando cada equa\c{c}\~ao \`a mesa real: a planta
$P(s)=A/s^2$, os ganhos $K_p,K_i,K_d$, a cinem\'atica inversa 3RPS e o
firmware em C rodando a 100\,Hz no ESP32-S3.
\end{tcolorbox}
\end{titlepage}

%------------------------------------------------------------------ PREFACIO ---
\chapter*{Como usar esta apostila}
\addcontentsline{toc}{chapter}{Como usar esta apostila}
\markboth{Como usar esta apostila}{}

Esta apostila foi escrita para quem est\'a construindo (e quer \emph{dominar}) o
\textbf{Ball Balancer 3RPS} --- uma mesa que equilibra uma bola no centro,
inclinando-se por tr\^es motores de passo sob comando de um controlador PID
digital embarcado. Cada cap\'itulo abre com a teoria geral e termina amarrando-a
ao \textbf{seu} sistema: os n\'umeros reais, o c\'odigo real e as decis\~oes de
projeto reais.

\begin{conceito}[Tr\^es tipos de caixa para guiar a leitura]
\textbf{Caixas azuis} (Conceito) fixam a defini\c{c}\~ao formal;
\textbf{caixas verdes} (No projeto) mostram onde aquilo aparece no Ball Balancer;
\textbf{caixas laranjas} (Exemplo num\'erico) fazem a conta com os par\^ametros
reais; \textbf{caixas roxas} (Intui\c{c}\~ao) traduzem a matem\'atica em imagem
mental; \textbf{caixas vermelhas} (Aten\c{c}\~ao) alertam para armadilhas.
\end{conceito}

\noindent\textbf{Pr\'e-requisitos:} c\'alculo (derivada e integral), no\c{c}\~oes
de equa\c{c}\~oes diferenciais e \'algebra linear b\'asica. N\~ao \'e preciso ter
visto Controle antes --- construiremos tudo do zero.

\tableofcontents
\listoffigures
\listoftables

%------------------------------------------------------- LISTA DE SIMBOLOS -----
\chapter*{Lista de S\'imbolos}
\addcontentsline{toc}{chapter}{Lista de S\'imbolos}
\markboth{Lista de S\'imbolos}{}

\begingroup
\renewcommand{\arraystretch}{1.3}
\begin{longtable}{>{\bfseries}l >{\RaggedRight\arraybackslash}p{10.5cm} l}
\toprule
S\'imbolo & Significado & Valor t\'ipico \\
\midrule\endhead
$x,y$        & Posi\c{c}\~ao da bola na mesa (eixos X e Y) & mm \\
$r$          & Refer\^encia / setpoint (centro da mesa) & $(93{,}5;\,70{,}5)$ mm \\
$e$          & Erro de posi\c{c}\~ao, $e = x - r$ & mm \\
$u,\;\alpha$ & Sa\'ida do controlador = inclina\c{c}\~ao da mesa & rad \\
$g$          & Acelera\c{c}\~ao da gravidade & $9810$ mm/s$^2$ \\
$A$          & Ganho da planta, $A=\tfrac{5}{7}g$ & $\approx 7007$ mm/s$^2$ \\
$P(s)$       & Fun\c{c}\~ao de transfer\^encia da planta & $A/s^2$ \\
$C(s)$       & Fun\c{c}\~ao de transfer\^encia do controlador (PID) & --- \\
$K_p,K_i,K_d$& Ganhos proporcional, integral e derivativo & ver Cap.~\ref{cap:ganhos} \\
$\zeta$      & Coeficiente de amortecimento & $0\!-\!\infty$ \\
$\omega_n$   & Frequ\^encia natural n\~ao amortecida & rad/s \\
$\omega_d$   & Frequ\^encia natural amortecida, $\omega_n\sqrt{1-\zeta^2}$ & rad/s \\
$M_p$        & Sobressinal (overshoot) m\'aximo & \% \\
$t_s,t_r,t_p$& Tempos de acomoda\c{c}\~ao, subida e de pico & s \\
$T$          & Per\'iodo de amostragem & $0{,}01$ s (100 Hz) \\
$\tau_d$     & Constante do filtro derivativo & $0{,}03$ s \\
$z$          & Vari\'avel da Transformada Z & --- \\
$\theta_A,\theta_B,\theta_C$ & \^Angulos dos tr\^es motores (cinem\'atica 3RPS) & graus \\
$(n_x,n_y)$  & Componentes do vetor normal da mesa & --- \\
\bottomrule
\end{longtable}
\endgroup

\mainmatter

%##############################################################################
\part{Fundamentos F\'isicos e Modelagem}
%##############################################################################

%==============================================================================
\chapter{Introdu\c{c}\~ao ao Problema}\label{cap:intro}
%==============================================================================
\index{Ball Balancer}

\section{O que \'e um Ball Balancer}
Um \textbf{Ball Balancer} (mesa equilibrista) \'e um sistema cujo \'unico objetivo
\'e manter uma \textbf{bola} parada num ponto de uma \textbf{superf\'icie} que pode
ser \textbf{inclinada}. Quando a mesa se inclina, a componente da gravidade ao
longo do tampo acelera a bola; o controlador deve inclinar a mesa \emph{na medida
certa} para frear a bola e traz\^e-la de volta ao centro.

\begin{projeto}
No nosso rig, a superf\'icie \'e uma \textbf{tela resistiva de 4 fios}
($187\times141$ mm) que mede a posi\c{c}\~ao da bola; tr\^es \textbf{motores de
passo NEMA 17} em arranjo \textbf{3RPS} inclinam a plataforma; e um \textbf{ESP32-S3}
roda dois controladores PID (um para X, outro para Y) a 100\,Hz.
\end{projeto}

\begin{figure}[h]\centering
\begin{tikzpicture}[scale=1.2]
  \fill[corprim!8] (-3,-1.4) rectangle (3,-1.0);
  \node[corcinza,font=\footnotesize] at (0,-1.2) {base / estrutura};
  \draw[corprim,very thick,rotate=10] (-2.5,0) -- (2.5,0);
  \shade[ball color=coracc,rotate=10] (1.1,0.32) circle (0.3);
  \draw[seta,rotate=10] (1.1,0.6) -- ++(0.9,0) node[right,font=\footnotesize]{a bola escorrega};
  \draw[seta,coralert] (1.55,0.0) -- ++(0,-0.8) node[below,font=\footnotesize]{$mg$};
  \draw[seta,corsec] (-2.5,-0.45) arc (200:240:2.2) node[midway,below,font=\footnotesize]{motor inclina};
  \node[font=\small\bfseries,corprim] at (0,1.3) {Bola sobre plano inclin\'avel};
\end{tikzpicture}
\caption{Princ\'ipio do Ball Balancer: inclinar o tampo gera a for\c{c}a que move a bola.}
\label{fig:principio}
\end{figure}

\section{Por que \'e um problema cl\'assico de controle}
O Ball Balancer (e seus primos \emph{ball-and-beam}, \emph{ball-on-plate}) \'e um
laborat\'orio did\'atico perfeito porque re\'une, num \'unico sistema barato e
vis\'ivel a olho nu:
\begin{itemize}[leftmargin=1.4em]
  \item uma planta \textbf{naturalmente inst\'avel} (a bola foge sozinha);
  \item din\^amica de \textbf{segunda ordem ``pura''} (duplo integrador), ideal
  para enxergar $\zeta$ e $\omega_n$;
  \item necessidade clara dos tr\^es termos \textbf{P}, \textbf{I} e \textbf{D};
  \item realimenta\c{c}\~ao por \textbf{sensor real} (ru\'ido, atraso, quantiza\c{c}\~ao);
  \item \textbf{atuadores reais} (motores com limite de velocidade e satura\c{c}\~ao).
\end{itemize}

\section{Por que \'e naturalmente inst\'avel}
\index{instabilidade}
Deixe a mesa parada e levemente torta: a bola \textbf{acelera} e \textbf{nunca
para por conta pr\'opria} --- ela rola at\'e cair. N\~ao existe ``for\c{c}a de
mola'' que a traga de volta. Em linguagem de controle, a planta tem
\textbf{dois polos na origem} (Cap.~\ref{cap:ft}): qualquer perturba\c{c}\~ao cresce sem
limite. \'E o controlador que \emph{cria}, artificialmente, a rigidez (via $K_p$)
e o amortecimento (via $K_d$) que a f\'isica n\~ao oferece.

\section{Os desafios do controle}
\begin{enumerate}[leftmargin=1.5em]
  \item \textbf{Estabilizar} algo que sozinho diverge (Cap.~\ref{cap:estab}).
  \item \textbf{Amortecer} sem deixar lento demais --- equil\'ibrio entre $K_p$ e $K_d$.
  \item \textbf{Eliminar o desvio} causado pela mesa torta --- papel do $K_i$ (e do
  \emph{feedforward} TRIM, Cap.~\ref{cap:futuro}).
  \item \textbf{Rejeitar ru\'ido} do sensor sem amplific\'a-lo no derivativo (Cap.~\ref{cap:filtro}).
  \item \textbf{Respeitar limites} f\'isicos: satura\c{c}\~ao de inclina\c{c}\~ao
  ($\pm0{,}25$ rad) e de velocidade dos motores (anti-windup, Cap.~\ref{cap:windup}).
  \item Fazer tudo isso em \textbf{tempo discreto}, a cada 10\,ms (Parte~III).
\end{enumerate}

%==============================================================================
\chapter{Revis\~ao Matem\'atica Compacta}\label{cap:matrev}
%==============================================================================
\index{revisao matematica@revis\~ao matem\'atica}

Antes de modelar a bola, vale alinhar cinco ferramentas. A meta aqui n\~ao \'e
rigor, e sim \textbf{intui\c{c}\~ao}: o que cada conceito \emph{significa} para o
nosso sistema.

\section{Derivada --- taxa de varia\c{c}\~ao}
A derivada mede \textbf{qu\~ao r\'apido} uma grandeza muda. Se $x(t)$ \'e a
posi\c{c}\~ao da bola,
\[
\dot x = \frac{dx}{dt}\ \text{(velocidade)},\qquad
\ddot x = \frac{d^2x}{dt^2}\ \text{(acelera\c{c}\~ao)}.
\]
\begin{intuicao}
Derivar \'e perguntar \emph{``est\'a aumentando ou diminuindo, e com que pressa?''}.
A velocidade \'e a derivada da posi\c{c}\~ao; a acelera\c{c}\~ao, a derivada da
velocidade. O termo \textbf{D} do PID usa exatamente isto: olha a \emph{velocidade}
da bola para antecipar para onde ela vai.
\end{intuicao}

\section{Integral --- acumula\c{c}\~ao}
A integral \'e o \textbf{oposto} da derivada: \textbf{acumula} (soma) a grandeza ao
longo do tempo --- a ``\'area sob a curva''.
\[
x(t) = x_0 + \int_0^t \dot x(\tau)\,d\tau.
\]
\begin{intuicao}
Sabendo a velocidade a cada instante e \emph{somando}, descobre-se o quanto a bola
andou. O termo \textbf{I} do PID acumula o \emph{erro} no tempo --- \'e a
``mem\'oria'' que percebe um desvio teimoso (mesa torta) e o corrige.
\end{intuicao}

\section{Equa\c{c}\~oes diferenciais}
Uma \textbf{equa\c{c}\~ao diferencial} relaciona uma fun\c{c}\~ao com suas
derivadas. A f\'isica da bola \'e uma delas:
\[
\ddot x = A\,\alpha \quad(\text{``a acelera\c{c}\~ao \'e proporcional \`a inclina\c{c}\~ao''}).
\]
Resolv\^e-la \'e achar $x(t)$ a partir dessa regra --- tarefa que a Transformada de
Laplace (Cap.~\ref{cap:laplace}) cumpre de forma \emph{alg\'ebrica}.

\section{Vari\'aveis de estado}
O \textbf{estado} \'e o conjunto m\'inimo de n\'umeros que resume tudo o que o
sistema precisa de si para evoluir. Para a bola num eixo,
\[
\mathbf{x} = \begin{bmatrix} \text{posi\c{c}\~ao} \\[2pt] \text{velocidade}\end{bmatrix}.
\]
\begin{intuicao}
Sabendo \emph{onde} a bola est\'a, \emph{com que velocidade}, e a inclina\c{c}\~ao
aplicada, o futuro fica determinado. ``Os dois n\'umeros que importam'' s\~ao
posi\c{c}\~ao e velocidade --- por isso a planta \'e de \textbf{segunda ordem}.
\end{intuicao}

\section{Sistema din\^amico}
Um \textbf{sistema din\^amico} \'e algo cujo estado \textbf{evolui no tempo} segundo
uma regra, em resposta a uma \textbf{entrada}:
\begin{center}
entrada (inclina\c{c}\~ao $\alpha$) $\;\rightarrow\;$ \fbox{sistema (bola)} $\;\rightarrow\;$ sa\'ida (posi\c{c}\~ao $x$).
\end{center}
\begin{projeto}
Todo este cap\'itulo serve a uma pergunta: \emph{``dada a inclina\c{c}\~ao que
comando, para onde vai a bola --- e como escolher a inclina\c{c}\~ao para que ela
fique no centro?''}. Derivada, integral e equa\c{c}\~ao diferencial s\~ao a
linguagem; o PID \'e a resposta.
\end{projeto}

%==============================================================================
\chapter{Modelagem F\'isica da Esfera}\label{cap:modelagem}
%==============================================================================
\index{modelagem}

Vamos deduzir, do zero, a equa\c{c}\~ao que governa a bola. O objetivo \'e chegar a
\begin{equation}
\boxed{\ \ddot{x} = \tfrac{5}{7}\,g\,\sin\alpha\ }
\label{eq:planta-nl}
\end{equation}
e entender \emph{de onde vem} o fator $5/7$.

\section{Cen\'ario e for\c{c}as atuantes}
Considere uma esfera maci\c{c}a de massa $m$ e raio $\rho$ que \textbf{rola sem
escorregar} sobre um plano inclinado de \^angulo $\alpha$. Ao longo do tampo agem:
\begin{itemize}[leftmargin=1.4em]
  \item a componente do peso $mg\sin\alpha$ (puxa a bola ladeira abaixo);
  \item a for\c{c}a de atrito est\'atico $F_{at}$ (no ponto de contato, \'e ela que
  faz a bola \emph{girar} em vez de deslizar).
\end{itemize}

\begin{figure}[h]\centering
\begin{tikzpicture}[scale=1.5]
  \draw[corprim,very thick] (-2.2,-0.9) -- (2.2,0.9);
  \shade[ball color=coracc] (0,0.42) circle (0.42);
  \draw[seta,coralert] (0,0) -- (0,-1.0) node[right,font=\footnotesize]{$mg$};
  \draw[seta,corsec] (0,0) -- (0.93,0.46) node[above,font=\footnotesize]{$mg\sin\alpha$};
  \draw[seta,corprim] ($(0,0)+(0.22,-0.42)$) -- ++(-0.6,-0.3) node[below,font=\footnotesize]{$F_{at}$};
  \draw[corcinza] (1.6,0.64) arc (27:0:1.0) ;
  \node[corcinza,font=\footnotesize] at (1.85,0.35) {$\alpha$};
\end{tikzpicture}
\caption{Diagrama de corpo livre da esfera rolando no plano inclinado.}
\label{fig:dcl}
\end{figure}

\section{Segunda Lei de Newton (transla\c{c}\~ao)}
Somando as for\c{c}as ao longo do tampo (sentido positivo ladeira abaixo):
\begin{equation}
m\,\ddot{x} = mg\sin\alpha - F_{at}.
\label{eq:newton}
\end{equation}

\section{Equa\c{c}\~ao de Euler (rota\c{c}\~ao)}
A \'unica for\c{c}a com bra\c{c}o de alavanca em torno do centro da bola \'e o
atrito (bra\c{c}o $=\rho$). Pela 2.\textsuperscript{a} lei de Newton para a
rota\c{c}\~ao,
\begin{equation}
I\,\dot{\omega} = F_{at}\,\rho,
\label{eq:euler}
\end{equation}
onde $I$ \'e o momento de in\'ercia e $\omega$ a velocidade angular.

\section{Rolamento sem escorregamento}
\index{rolamento sem escorregamento}
A condi\c{c}\~ao de rolar sem deslizar amarra transla\c{c}\~ao e rota\c{c}\~ao:
\begin{equation}
\dot{x} = \omega\,\rho \quad\Longrightarrow\quad \ddot{x} = \dot{\omega}\,\rho
\quad\Longrightarrow\quad \dot{\omega} = \frac{\ddot{x}}{\rho}.
\label{eq:roll}
\end{equation}

\section{Eliminando o atrito: a origem do 5/7}
Para a \textbf{esfera maci\c{c}a}, o momento de in\'ercia em torno do centro \'e
\begin{equation}
I = \tfrac{2}{5}\,m\rho^2.
\end{equation}
Substituindo \eqref{eq:roll} em \eqref{eq:euler}:
\begin{equation}
F_{at} = \frac{I\,\dot\omega}{\rho}
       = \frac{1}{\rho}\Big(\tfrac{2}{5}m\rho^2\Big)\frac{\ddot{x}}{\rho}
       = \tfrac{2}{5}\,m\,\ddot{x}.
\end{equation}
Levando esse $F_{at}$ de volta \`a equa\c{c}\~ao de Newton \eqref{eq:newton}:
\begin{equation}
m\ddot{x} = mg\sin\alpha - \tfrac{2}{5}m\ddot{x}
\;\Longrightarrow\;
\Big(1+\tfrac{2}{5}\Big)\ddot{x} = g\sin\alpha
\;\Longrightarrow\;
\tfrac{7}{5}\ddot{x} = g\sin\alpha,
\end{equation}
de onde sai, finalmente, a Eq.~\eqref{eq:planta-nl}: $\ddot{x} = \tfrac{5}{7}g\sin\alpha$.

\begin{intuicao}[Por que ``apenas'' $5/7$ da gravidade?]
Parte da energia que a gravidade fornece n\~ao vira velocidade de
\emph{transla\c{c}\~ao}: ela vai \textbf{girar} a bola. Uma esfera maci\c{c}a
``gasta'' $2/7$ do empurr\~ao para rodar e sobra $5/7$ para andar. Por isso uma
bola desce a rampa mais devagar do que um bloco que desliza sem girar (esse teria
fator $1$). Um anel/casca esf\'erica desceria ainda mais devagar.
\end{intuicao}

\section{Comparando esfera, cilindro e anel}
O fator depende s\'o de \emph{como a massa se distribui} (do momento de in\'ercia
$I=c\,m\rho^2$). Refazendo a dedu\c{c}\~ao com um $I$ gen\'erico, a acelera\c{c}\~ao
fica $\ddot x=\dfrac{g\sin\alpha}{1+c}$:

\begin{table}[h]\centering
\caption{Quanto da gravidade vira ``andar'' para cada corpo que rola.}
\label{tab:corpos}
\begin{tabular}{lccc}
\toprule
Corpo & $I$ & fator $\dfrac{1}{1+c}$ & desce\dots \\
\midrule
Bloco deslizando (n\~ao gira) & $0$ & $1$ & mais r\'apido \\
\textbf{Esfera maci\c{c}a (nossa bola)} & $\tfrac{2}{5}m\rho^2$ & $\mathbf{5/7\approx0{,}71}$ & --- \\
Cilindro maci\c{c}o & $\tfrac{1}{2}m\rho^2$ & $2/3\approx0{,}67$ & um pouco mais devagar \\
Anel / casca fina & $m\rho^2$ & $1/2=0{,}50$ & bem mais devagar \\
\bottomrule
\end{tabular}
\end{table}

\begin{intuicao}[Para onde vai a energia da gravidade]
A energia potencial perdida pela bola ao descer vira energia cin\'etica, que se
\textbf{divide} em duas partes:
\[
\underbrace{mg\,h}_{\text{potencial}} \;=\;
\underbrace{\tfrac12 m v^2}_{\text{transla\c{c}\~ao (andar)}} +
\underbrace{\tfrac12 I \omega^2}_{\text{rota\c{c}\~ao (girar)}}.
\]
Quanto mais massa longe do centro (maior $I$), mais energia ``escapa'' para girar e
menos sobra para andar. A esfera maci\c{c}a concentra massa perto do centro, ent\~ao
guarda boa parte (5/7) para a transla\c{c}\~ao --- por isso \'e o corpo ideal para um
Ball Balancer \'agil.
\end{intuicao}

\begin{projeto}[O fator vive no c\'odigo]
Na simula\c{c}\~ao f\'isica (\texttt{tools/PIDSimba.py}) esse mesmo coeficiente
aparece como \texttt{ROLL\_FACTOR = 5.0/7.0}. A f\'isica do firmware e a da
simula\c{c}\~ao usam exatamente a Eq.~\eqref{eq:planta-nl}.
\end{projeto}

%==============================================================================
\chapter{Lineariza\c{c}\~ao}\label{cap:linear}
%==============================================================================
\index{linearizacao@lineariza\c{c}\~ao}

A Eq.~\eqref{eq:planta-nl} \'e \textbf{n\~ao linear} por causa do $\sin\alpha$.
Ferramentas como Laplace, polos e $\zeta/\omega_n$ exigem um modelo \textbf{linear}.
A sa\'ida \'e a aproxima\c{c}\~ao de pequenos \^angulos.

\section{A aproxima\c{c}\~ao $\sin\alpha \approx \alpha$}
A s\'erie de Taylor do seno em torno de zero \'e
\begin{equation}
\sin\alpha = \alpha - \frac{\alpha^3}{6} + \frac{\alpha^5}{120} - \cdots
\end{equation}
Para $\alpha$ pequeno (em \textbf{radianos}), o termo dominante \'e o pr\'oprio
$\alpha$, e os demais s\~ao desprez\'iveis. Logo $\sin\alpha\approx\alpha$ e a
planta vira \textbf{linear}:
\begin{equation}
\boxed{\ \ddot{x} \approx \tfrac{5}{7}\,g\,\alpha \ } .
\label{eq:planta-lin}
\end{equation}

\section{Quando a aproxima\c{c}\~ao \'e v\'alida}
\begin{exemplo}[Erro da aproxima\c{c}\~ao no limite de opera\c{c}\~ao]
O firmware satura a inclina\c{c}\~ao em $|\alpha|\le 0{,}25$ rad ($\approx14{,}3^\circ$).
No pior caso:
\[
\sin(0{,}25)=0{,}24740,\qquad \text{erro relativo} = \frac{0{,}25-0{,}24740}{0{,}24740}
\approx 1{,}05\%.
\]
Ou seja, mesmo no \textbf{extremo} do curso a lineariza\c{c}\~ao erra cerca de
$1\%$ --- e na opera\c{c}\~ao normal (poucos graus) o erro \'e bem menor que isso.
\end{exemplo}

\begin{table}[h]\centering
\caption{Qualidade da aproxima\c{c}\~ao $\sin\alpha\approx\alpha$.}
\label{tab:lin}
\begin{tabular}{cccc}
\toprule
$\alpha$ (graus) & $\alpha$ (rad) & $\sin\alpha$ & erro relativo \\
\midrule
$1^\circ$  & 0,0175 & 0,01745 & 0,005\% \\
$5^\circ$  & 0,0873 & 0,08716 & 0,13\% \\
$10^\circ$ & 0,1745 & 0,17365 & 0,51\% \\
$14{,}3^\circ$ & 0,2497 & 0,24740 & 1,05\% \\
\bottomrule
\end{tabular}
\end{table}

\section{O que acontece se ela n\~ao valer}
\begin{atencao}[Limita\c{c}\~ao consciente do projeto]
Para \^angulos grandes, $\sin\alpha < \alpha$: a mesa entregaria \emph{menos}
acelera\c{c}\~ao do que o modelo linear sup\~oe, e o controlador projetado ficaria
``otimista'' (ganho efetivo menor que o previsto). Por isso a satura\c{c}\~ao em
$0{,}25$ rad n\~ao \'e s\'o seguran\c{c}a mec\^anica: ela \textbf{mant\'em o
sistema dentro da regi\~ao onde a teoria linear vale}. Operar a mesa em \^angulos
muito maiores invalidaria todo o projeto de ganhos da Parte~II.
\end{atencao}

%==============================================================================
\chapter{Transformada de Laplace}\label{cap:laplace}
%==============================================================================
\index{Laplace}

\section{O que \'e e por que usamos}
A \textbf{Transformada de Laplace} converte uma fun\c{c}\~ao do tempo $f(t)$ numa
fun\c{c}\~ao de uma vari\'avel complexa $s$:
\begin{equation}
F(s) = \mathcal{L}\{f(t)\} = \int_0^{\infty} f(t)\,e^{-st}\,dt.
\end{equation}
A m\'agica: ela transforma \textbf{derivadas e integrais} (opera\c{c}\~oes do
c\'alculo) em \textbf{multiplica\c{c}\~oes e divis\~oes} por $s$ (opera\c{c}\~oes de
\'algebra). Uma equa\c{c}\~ao diferencial vira uma equa\c{c}\~ao alg\'ebrica.

\begin{conceito}[As duas propriedades que usaremos o tempo todo]
Com condi\c{c}\~oes iniciais nulas,
\[
\mathcal{L}\{\dot f(t)\} = s\,F(s), \qquad
\mathcal{L}\!\left\{\textstyle\int_0^t f\,d\tau\right\} = \frac{F(s)}{s}.
\]
\textbf{Derivar} $\leftrightarrow$ multiplicar por $s$. \textbf{Integrar}
$\leftrightarrow$ dividir por $s$.
\end{conceito}

\section{Derivando $P(s)=A/s^2$ passo a passo}
Partimos da planta linearizada \eqref{eq:planta-lin}, chamando a entrada de
$u=\alpha$ (a inclina\c{c}\~ao comandada) e $A=\tfrac{5}{7}g$:
\begin{align}
\ddot{x}(t) &= A\,u(t). &&\text{(equa\c{c}\~ao no tempo)}\\
\mathcal{L}\{\ddot{x}(t)\} &= \mathcal{L}\{A\,u(t)\}. &&\text{(aplica Laplace)}\\
s^2 X(s) &= A\,U(s). &&\text{(derivada segunda }\to s^2)\\
\frac{X(s)}{U(s)} &= \frac{A}{s^2}. &&\text{(isola a raz\~ao sa\'ida/entrada)}
\end{align}
\begin{equation}
\boxed{\ P(s)=\dfrac{X(s)}{U(s)}=\dfrac{A}{s^{2}},\qquad A\approx 7007\ \text{mm/s}^2\ }
\label{eq:plant}
\end{equation}

\begin{projeto}[De onde sai o $A\approx 7007$]
$A=\tfrac{5}{7}\,g$ com $g=9810$ mm/s$^2$ (gravidade em mil\'imetros, pois medimos
a bola em mm). Logo $A=\tfrac{5}{7}\cdot 9810 \approx 7007$ mm/s$^2$ por radiano de
inclina\c{c}\~ao. Esse n\'umero \'e a ``alavanca'' que liga inclina\c{c}\~ao a
acelera\c{c}\~ao em todo o projeto.
\end{projeto}

%==============================================================================
\chapter{Fun\c{c}\~ao de Transfer\^encia da Planta}\label{cap:ft}
%==============================================================================
\index{funcao de transferencia@fun\c{c}\~ao de transfer\^encia}\index{planta}

\section{O que \'e uma ``planta''}
\textbf{Planta} \'e o sistema f\'isico que queremos controlar --- aqui, ``a bola
respondendo \`a inclina\c{c}\~ao''. Sua fun\c{c}\~ao de transfer\^encia
$P(s)=A/s^2$ resume \emph{toda} a din\^amica: dada uma inclina\c{c}\~ao $U(s)$, ela
diz qual ser\'a a posi\c{c}\~ao $X(s)$.

\begin{intuicao}[Lendo $P(s)=A/s^2$ por partes]
\textbf{Numerador ($A$):} \'e o ``ganho/alavanca'' da planta --- \emph{quanta}
acelera\c{c}\~ao cada radiano de inclina\c{c}\~ao gera ($A\approx7007$ mm/s$^2$).
Como n\~ao h\'a $s$ no numerador, \textbf{n\~ao h\'a zeros}: a planta n\~ao
``real\c{c}a'' nenhuma frequ\^encia, ela apenas integra.\\
\textbf{Denominador ($s^2$):} \'e a \textbf{din\^amica}. Cada $s$ no denominador
significa ``integrar uma vez''; dois $s$ significam integrar \emph{duas} vezes.
\end{intuicao}

\begin{intuicao}[Por que surgem \emph{dois} integradores --- e de forma natural]
\'E pura f\'isica, n\~ao escolha de projeto. A inclina\c{c}\~ao gera
\textbf{acelera\c{c}\~ao} ($\ddot x = A\alpha$); mas o que o sensor mede \'e
\textbf{posi\c{c}\~ao}. Para ir de uma \`a outra, a natureza \emph{integra duas
vezes}: acelera\c{c}\~ao $\to$ velocidade $\to$ posi\c{c}\~ao. Esses dois ``integrar''
s\~ao, exatamente, os dois fatores $1/s$ --- e os dois polos na origem.
\end{intuicao}

\section{Polos e zeros: significado f\'isico}
\index{polos}\index{zeros}
\begin{conceito}[Defini\c{c}\~ao]
\textbf{Polos} s\~ao as ra\'izes do \emph{denominador}; ditam a \textbf{forma
natural} da resposta (cresce? decai? oscila?). \textbf{Zeros} s\~ao as ra\'izes do
\emph{numerador}; moldam \emph{como} os modos s\~ao excitados. Em $P(s)=A/s^2$ o
numerador \'e a constante $A$: \textbf{n\~ao h\'a zeros}, e h\'a um \textbf{polo
duplo em} $s=0$.
\end{conceito}

\begin{figure}[h]\centering
\begin{tikzpicture}[scale=1.0]
  \draw[seta] (-2.6,0) -- (2.6,0) node[right,font=\footnotesize]{$\mathrm{Re}\,s$};
  \draw[seta] (0,-2.0) -- (0,2.0) node[above,font=\footnotesize]{$\mathrm{Im}\,s$};
  \node[coralert] at (0,0) {\large $\boldsymbol\times$};
  \node[coralert] at (0.18,0.22) {\large $\boldsymbol\times$};
  \node[corcinza,font=\footnotesize,anchor=west] at (0.35,-0.35) {polo duplo em $s=0$};
  \draw[decorate,decoration={brace,amplitude=4pt},corsec] (-2.4,-1.6) -- (-0.1,-1.6)
     node[midway,below,font=\footnotesize,corsec]{semiplano est\'avel};
  \draw[decorate,decoration={brace,amplitude=4pt},coralert] (0.1,1.6) -- (2.4,1.6)
     node[midway,above,font=\footnotesize,coralert]{semiplano inst\'avel};
\end{tikzpicture}
\caption{Plano-$s$ da planta: dois polos exatamente na origem (fronteira da estabilidade).}
\label{fig:splane}
\end{figure}

\begin{intuicao}[Polo = a ``tend\^encia natural'' do sistema]
Um polo descreve como o sistema se comporta \emph{sozinho}, sem ningu\'em mexer ---
e a posi\c{c}\~ao dele no plano-$s$ j\'a conta a hist\'oria toda.
\end{intuicao}

\begin{table}[h]\centering
\caption{Onde est\'a o polo $\Rightarrow$ como o sistema reage (com o exemplo da bola).}
\label{tab:polos}
\small
\begin{tabularx}{\textwidth}{l l X}
\toprule
Polo\dots & Comportamento natural & No Ball Balancer \\
\midrule
\`a \textbf{esquerda} ($\sigma<0$) & decai sozinho (\textbf{est\'avel}) & o alvo: o controlador puxa os polos para c\'a \\
na \textbf{origem} ($\sigma=0$) & nem cresce nem decai (\textbf{marginal}) & a planta crua --- a bola n\~ao volta nem cai sozinha \\
\`a \textbf{direita} ($\sigma>0$) & cresce sem limite (\textbf{inst\'avel}) & sinal/ganho errado: a mesa \emph{empurra} a bola para fora \\
\bottomrule
\end{tabularx}
\end{table}

\section{O que significam dois polos na origem}
Cada $1/s$ \'e um \textbf{integrador}. Dois polos na origem $=$ \textbf{duplo
integrador}: a entrada \'e integrada \emph{duas vezes} at\'e virar posi\c{c}\~ao.
\begin{equation}
u \;\xrightarrow{\ \int\ }\; \dot{x}\;(\text{velocidade}) \;\xrightarrow{\ \int\ }\; x\;(\text{posi\c{c}\~ao}).
\end{equation}

\begin{intuicao}[Por que a bola acelera ``para sempre''?]
Incline a mesa por um \^angulo fixo $\alpha_0$. A acelera\c{c}\~ao $\ddot x=A\alpha_0$
\'e constante e n\~ao nula. Integrando uma vez, a \textbf{velocidade cresce
linearmente} ($\dot x = A\alpha_0 t$); integrando de novo, a \textbf{posi\c{c}\~ao
cresce com o quadrado do tempo} ($x=\tfrac12 A\alpha_0 t^2$). Nada nessa cadeia
``puxa a bola de volta'': sem realimenta\c{c}\~ao, ela diverge. Esse \'e o
significado f\'isico, palp\'avel, dos dois polos na origem.
\end{intuicao}

\begin{conceito}[Sistema do Tipo 2]
Ter dois integradores torna a planta do \textbf{Tipo 2}. Consequ\^encia
positiva: em malha fechada, o erro estacion\'ario a uma \textbf{refer\^encia
constante} e at\'e a uma \textbf{rampa} \'e \emph{nulo} --- antes mesmo de
adicionar o termo integral do PID. O $K_i$, aqui, serve principalmente contra
\textbf{perturba\c{c}\~oes} (a mesa torta), n\~ao contra erro de refer\^encia.
\end{conceito}

%##############################################################################
\part{An\'alise de Sistemas de Controle}
%##############################################################################

%==============================================================================
\chapter{Diagramas de Blocos}\label{cap:blocos}
%==============================================================================
\index{diagrama de blocos}

Diagramas de blocos mostram o \textbf{fluxo de sinais}. Vamos do mais detalhado
(o sistema f\'isico real) ao mais abstrato (malha fechada cl\'assica).

\section{Sistema real (cascata completa)}
\begin{figure}[h]\centering
\begin{tikzpicture}[node distance=11mm and 12mm]
  \node[soma] (s) {};
  \node[left=of s] (r) {$r$};
  \node[bloco,right=of s] (pid) {PID\\ $C(s)$};
  \node[bloco,right=of pid] (kin) {Cinem\'atica\\ inversa 3RPS};
  \node[bloco,right=of kin] (mot) {Motores\\ de passo};
  \node[blocoS,right=of mot] (plat) {Plataforma};
  \node[blocoS,right=of plat] (ball) {Bola\\ $A/s^2$};
  \draw[seta] (r) -- (s);
  \draw[seta] (s) -- node[rotulo,above]{$e$} (pid);
  \draw[seta] (pid) -- node[rotulo,above]{$n_x,n_y$} (kin);
  \draw[seta] (kin) -- node[rotulo,above]{$\theta_{A,B,C}$} (mot);
  \draw[seta] (mot) -- (plat);
  \draw[seta] (plat) -- node[rotulo,above]{$\alpha$} (ball);
  \draw[seta] (ball) -- ++(0,-1.4) -| node[rotulo,pos=0.25,below]{sensor (tela)} (s);
  \node[font=\footnotesize] at ($(s)+(-0.0,0.32)$) {$-$};
  \node[above=1mm of ball,font=\footnotesize] {$x,y$};
\end{tikzpicture}
\caption{Cadeia f\'isica completa: sensor $\to$ PID $\to$ cinem\'atica $\to$ motores $\to$ plataforma $\to$ bola.}
\label{fig:cascata}
\end{figure}

\section{Sistema simplificado (controle puro)}
A cinem\'atica e os motores s\~ao r\'apidos e essencialmente \textbf{est\'aticos}
perante a din\^amica lenta da bola (Cap.~\ref{cap:arq}): podemos coloc\'a-los como ganho
unit\'ario e enxergar o miolo do controle.
\begin{figure}[h]\centering
\begin{tikzpicture}[node distance=12mm and 16mm]
  \node[soma] (s) {};
  \node[left=of s] (r) {$r$};
  \node[bloco,right=of s] (pid) {$C(s)=K_p+\dfrac{K_i}{s}+K_d s$};
  \node[blocoS,right=of pid] (p) {$P(s)=\dfrac{A}{s^2}$};
  \node[right=16mm of p] (y) {$x$};
  \draw[seta] (r) -- (s);
  \draw[seta] (s) -- node[rotulo,above]{$e$} (pid);
  \draw[seta] (pid) -- node[rotulo,above]{$u$} (p);
  \draw[seta] (p) -- (y);
  \draw[seta] (p) -- ++(0,-1.3) -| (s);
  \node[font=\footnotesize] at ($(s)+(0,0.32)$) {$-$};
\end{tikzpicture}
\caption{Malha de controle reduzida ao essencial: controlador $C(s)$ sobre a planta $P(s)$.}
\label{fig:reduzido}
\end{figure}

\section{Malha fechada e realimenta\c{c}\~ao negativa}
\index{realimentacao negativa@realimenta\c{c}\~ao negativa}
A fun\c{c}\~ao de transfer\^encia de malha fechada (entrada $r$, sa\'ida $x$) \'e
\begin{equation}
\frac{X(s)}{R(s)} = \frac{C(s)P(s)}{1+C(s)P(s)}.
\label{eq:mf}
\end{equation}
O sinal de \textbf{menos} no somador (realimenta\c{c}\~ao \emph{negativa}) \'e o que
faz o sistema \emph{corrigir} o erro em vez de amplific\'a-lo.

\begin{projeto}[O sinal no firmware: \texttt{TILT\_SIGN} $=-1$]
O PID calcula $u=K_p\,e+\cdots$ com $e=\text{medida}-\text{setpoint}$. Para a mesa
inclinar no sentido que traz a bola de volta, aplica-se $n_x=-u$
(\texttt{TILT\_SIGN\_X = -1.0f} em \texttt{control.c}). Se um eixo empurrar a bola
para longe, basta inverter aquele sinal ao vivo: \texttt{SX +1} ou \texttt{SY +1}.
\'E a Eq.~\eqref{eq:mf} ganhando vida no hardware.
\end{projeto}

%==============================================================================
\chapter{Estabilidade}\label{cap:estab}
%==============================================================================
\index{estabilidade}

\section{Est\'avel, inst\'avel, marginalmente est\'avel}
\begin{conceito}[Crit\'erio pelos polos de malha fechada]
Um sistema linear \'e \textbf{est\'avel} se \emph{todos} os polos de malha fechada
t\^em parte real \textbf{negativa} (semiplano esquerdo); \textbf{inst\'avel} se
algum tem parte real \textbf{positiva}; \textbf{marginalmente est\'avel} se h\'a
polos \emph{exatamente} sobre o eixo imagin\'ario (e nenhum \`a direita).
\end{conceito}

\section{Onde a planta se encaixa}
A planta isolada $A/s^2$ tem polo duplo em $s=0$ --- \textbf{em cima} do eixo
imagin\'ario: \emph{marginalmente est\'avel} no melhor caso e, na pr\'atica
(qualquer perturba\c{c}\~ao), divergente. \textbf{Sem controlador, a mesa n\~ao
funciona.}

\section{Como o controlador estabiliza}
Feche a malha com um \textbf{PD} ($C(s)=K_p+K_d s$) sobre $A/s^2$. A
equa\c{c}\~ao caracter\'istica $1+C(s)P(s)=0$ vira
\begin{equation}
s^2 + A K_d\, s + A K_p = 0.
\label{eq:carac-pd}
\end{equation}
Pelo crit\'erio de Routh--Hurwitz, um polin\^omio de 2.\textsuperscript{a} ordem
$s^2+a_1 s+a_0$ tem ambas as ra\'izes no semiplano esquerdo \emph{se e somente se}
$a_1>0$ e $a_0>0$. Aqui $a_1=AK_d$ e $a_0=AK_p$: basta $K_p>0$ e $K_d>0$.
\begin{atencao}[O $K_d$ \'e quem traz a estabilidade]
Sem o termo derivativo ($K_d=0$), a Eq.~\eqref{eq:carac-pd} vira $s^2+AK_p=0$, com
ra\'izes $\pm j\sqrt{AK_p}$ --- \emph{de novo} sobre o eixo imagin\'ario:
oscila\c{c}\~ao n\~ao amortecida. \textbf{\'E o $K_d$ que empurra os polos para a
esquerda} e amortece. Guarde isto: num duplo integrador, derivativo n\~ao \'e
``tempero'', \'e requisito de estabilidade.
\end{atencao}

%==============================================================================
\chapter{O Controlador PID}\label{cap:pid}
%==============================================================================
\index{PID}

\section{A lei de controle}
\begin{equation}
u(t) = \underbrace{K_p\,e(t)}_{\text{presente}}
     + \underbrace{K_i\!\int_0^t\! e(\tau)\,d\tau}_{\text{passado}}
     + \underbrace{K_d\,\dot e(t)}_{\text{futuro}},
\qquad C(s)=K_p+\frac{K_i}{s}+K_d s.
\end{equation}
Cada termo olha para um ``tempo'': o \textbf{P} reage ao erro \emph{agora}, o
\textbf{I} acumula o \emph{passado}, e o \textbf{D} antecipa a tend\^encia
(\emph{futuro}).

\begin{intuicao}[As tr\^es analogias mec\^anicas]
A forma mais f\'acil de ``sentir'' o PID \'e imagin\'a-lo como tr\^es pe\c{c}as
mec\^anicas agindo sobre a bola:
\begin{itemize}[leftmargin=1.4em]
  \item \textbf{$K_p$ = mola virtual.} Puxa a bola de volta com for\c{c}a
  proporcional ao \emph{quanto} ela se afastou. S\'o mola, por\'em, oscila.
  \item \textbf{$K_d$ = amortecedor.} Op\~oe-se \`a \emph{velocidade} da bola, como
  um amortecedor de carro. \'E ele que ``segura'' a oscila\c{c}\~ao da mola.
  \item \textbf{$K_i$ = mem\'oria acumulativa.} Lembra de um desvio que
  \emph{insiste} (mesa torta) e vai inclinando mais at\'e zer\'a-lo.
\end{itemize}
Mola + amortecedor + mem\'oria: \'e literalmente um sistema massa-mola-amortecedor
que \emph{voc\^e} projeta via $K_p,K_d$ e ajusta o ponto de equil\'ibrio via $K_i$.
\end{intuicao}

\begin{table}[h]\centering
\caption{Efeito de aumentar cada ganho sobre a trajet\'oria da bola.}
\label{tab:pidresumo}
\small
\begin{tabularx}{\textwidth}{l l X X}
\toprule
Ganho & Analogia & Se \textbf{aumenta} & Se \textbf{falta} \\
\midrule
$K_p$ & mola & mais r\'apido; em excesso, \textbf{oscila} & mole, lento, n\~ao chega \\
$K_d$ & amortecedor & menos overshoot, mais firme & \textbf{oscila}/instabiliza; amplifica ru\'ido se em excesso \\
$K_i$ & mem\'oria & zera o desvio fixo & desvio fixo \emph{persiste}; em excesso, \emph{windup} \\
\bottomrule
\end{tabularx}
\end{table}

\section{$K_p$ --- proporcional}
\index{Kp@$K_p$}
\textbf{Interpreta\c{c}\~ao f\'isica:} rigidez de uma ``mola virtual''. Quanto mais
longe do centro, mais a mesa inclina para empurrar a bola de volta.
\begin{itemize}[leftmargin=1.4em]
  \item \textbf{Aumenta $K_p$}: resposta mais r\'apida ($\omega_n=\sqrt{AK_p}$
  sobe); em excesso, oscila e fica \`a beira da instabilidade.
  \item \textbf{Diminui $K_p$}: resposta mole e lenta; a bola demora a voltar.
\end{itemize}

\section{$K_d$ --- derivativo}
\index{Kd@$K_d$}
\textbf{Interpreta\c{c}\~ao f\'isica:} amortecimento que reage \`a \emph{velocidade}
da bola. Se a bola corre para o centro r\'apido demais, o D ``segura'' antes que ela
passe do ponto.
\begin{itemize}[leftmargin=1.4em]
  \item \textbf{Aumenta $K_d$}: mais amortecimento, menos overshoot; em excesso,
  lentid\~ao e \textbf{amplifica\c{c}\~ao de ru\'ido} (Cap.~\ref{cap:filtro}).
  \item \textbf{Diminui $K_d$}: mais overshoot e oscila\c{c}\~ao; em $K_d\to0$,
  instabilidade marginal (Cap.~\ref{cap:estab}).
\end{itemize}

\section{$K_i$ --- integral}
\index{Ki@$K_i$}
\textbf{Interpreta\c{c}\~ao f\'isica:} mem\'oria que elimina desvio
\emph{persistente}. Se a bola insiste em ficar deslocada (mesa torta), o integral
acumula esse erro e adiciona a inclina\c{c}\~ao constante que falta.
\begin{itemize}[leftmargin=1.4em]
  \item \textbf{Aumenta $K_i$}: zera o desvio mais r\'apido; em excesso, causa
  \emph{windup} e oscila\c{c}\~ao lenta.
  \item \textbf{Diminui $K_i$}: desvio some devagar (ou ``nunca'', na pr\'atica).
\end{itemize}

\begin{figure}[h]\centering
\begin{tikzpicture}
\begin{axis}[width=12.5cm,height=6.2cm,xlabel={tempo (s)},ylabel={posi\c{c}\~ao (norm.)},
  xmin=0,xmax=6,ymin=0,ymax=1.6,legend pos=south east,grid=both,
  grid style={corcinza!20},legend style={font=\footnotesize},
  title={Efeito de $K_d$ na resposta (mesmo $K_p$, $\omega_n=2{,}19$)},
  every axis title/.style={font=\small\bfseries,color=corprim}]
  \addplot[coralert,very thick,domain=0:6,samples=160]
    {1 - exp(-0.15*2.19*x)/sqrt(1-0.15^2)*sin(deg(2.19*sqrt(1-0.15^2)*x)+acos(0.15))};
  \addlegendentry{$K_d$ baixo ($\zeta=0{,}15$): oscila}
  \addplot[coracc,very thick,domain=0:6,samples=160]
    {1 - exp(-0.5*2.19*x)/sqrt(1-0.5^2)*sin(deg(2.19*sqrt(1-0.5^2)*x)+acos(0.5))};
  \addlegendentry{$K_d$ m\'edio ($\zeta=0{,}5$)}
  \addplot[corsec,very thick,domain=0:6,samples=160]
    {1-(1+2.19*x)*exp(-2.19*x)};
  \addlegendentry{$K_d$ alto ($\zeta=1$): sem overshoot}
  \draw[corcinza,dashed] (axis cs:0,1) -- (axis cs:6,1);
\end{axis}
\end{tikzpicture}
\caption{Mais derivativo $\Rightarrow$ mais amortecimento: a oscila\c{c}\~ao some \`a custa de um pouco de velocidade.}
\label{fig:kd}
\end{figure}

\begin{projeto}[Os tr\^es termos no seu firmware]
Em \texttt{pid.c}, a linha \texttt{out = p->kp*err + p->ki*p->integ + p->kd*deriv;}
\'e literalmente $u=K_p e+K_i\!\int e+K_d\dot e$. O $K_i$ combate a
\textbf{base torta} (Cap.~\ref{cap:futuro}), e o firmware ainda oferece o \textbf{TRIM}, um
\emph{feedforward} que faz o mesmo trabalho do integral, por\'em instantaneamente.
\end{projeto}

%==============================================================================
\chapter{Como os Ganhos Foram Obtidos}\label{cap:ganhos}
%==============================================================================
\index{alocacao de polos@aloca\c{c}\~ao de polos}

\section{Tr\^es caminhos para escolher $K_p,K_i,K_d$}
\begin{enumerate}[leftmargin=1.5em]
  \item \textbf{Aloca\c{c}\~ao de polos} (anal\'itico): escolhe-se \emph{onde} os
  polos de malha fechada devem ficar (a partir de requisitos de $\zeta,\omega_n$) e
  resolve-se para os ganhos. \'E o m\'etodo do TCC de refer\^encia~\cite{cordeiro}.
  \item \textbf{Ziegler--Nichols} (experimental): leva o sistema ao limite da
  oscila\c{c}\~ao e usa tabelas (Cap.~\ref{cap:zn}).
  \item \textbf{Tentativa e erro guiada} (pr\'atico): sobe $K_p$, depois $K_d$,
  depois um toque de $K_i$ --- \'e como se faz o ajuste fino no rig (Cap.~\ref{cap:sintonia}).
\end{enumerate}

\section{Aloca\c{c}\~ao de polos passo a passo}
A malha fechada do PID sobre $A/s^2$ \'e de \textbf{terceira ordem}. Sua
equa\c{c}\~ao caracter\'istica ($1+C(s)P(s)=0$, com $C=K_p+K_i/s+K_d s$) \'e
\begin{equation}
s^3 + A K_d\,s^2 + A K_p\,s + A K_i = 0.
\label{eq:carac3}
\end{equation}
Escolhemos um polin\^omio \textbf{desejado}: um par dominante de 2.\textsuperscript{a}
ordem (define $\zeta,\omega_n$) mais um polo real $p_o$ r\'apido:
\begin{equation}
(s^2 + 2\zeta\omega_n s + \omega_n^2)(s+p_o)
= s^3 + (2\zeta\omega_n+p_o)s^2 + (\omega_n^2+2\zeta\omega_n p_o)s + \omega_n^2 p_o.
\end{equation}
Igualando coeficiente a coeficiente com \eqref{eq:carac3}:
\begin{equation}
\boxed{\;K_d=\dfrac{2\zeta\omega_n+p_o}{A},\quad
        K_p=\dfrac{\omega_n^2+2\zeta\omega_n p_o}{A},\quad
        K_i=\dfrac{\omega_n^2 p_o}{A}\;}
\label{eq:ganhos}
\end{equation}

\subsection{Dos requisitos a $\zeta$ e $\omega_n$}
Requisitos do projeto: sobressinal $PO\le 7{,}5\%$ e acomoda\c{c}\~ao
$t_s\le 2{,}5$ s. Pelas f\'ormulas de 2.\textsuperscript{a} ordem (Cap.~\ref{cap:zeta}--\ref{cap:tempos}):
\begin{equation}
\zeta=\sqrt{\frac{\ln^2(PO/100)}{\pi^2+\ln^2(PO/100)}}\approx 0{,}636,
\qquad
\omega_n\approx 2{,}19\ \text{rad/s}.
\end{equation}

\begin{exemplo}[Ganhos resultantes (com $p_o=1$ e $A=7007$)]
Substituindo $\zeta=0{,}636$, $\omega_n=2{,}19$, $p_o=1$ em \eqref{eq:ganhos}:
\[
K_d=\frac{2(0{,}636)(2{,}19)+1}{7007}\approx 5{,}4\times10^{-4},\quad
K_p\approx 1{,}08\times10^{-3},\quad
K_i\approx 6{,}83\times10^{-4}.
\]
Esses s\~ao os ganhos ``de projeto'' usados na simula\c{c}\~ao
(\texttt{sim\_pid.py}/\texttt{PIDSimba.py}).
\end{exemplo}

\section{Comparativo dos conjuntos de ganhos do projeto}
\begin{table}[h]\centering
\caption{Conjuntos de ganhos: projeto (anal\'itico), default do firmware e o que est\'a em opera\c{c}\~ao.}
\label{tab:ganhos}
\small
\begin{tabularx}{\textwidth}{l c c c X}
\toprule
Conjunto & $K_p$ & $K_i$ & $K_d$ & Observa\c{c}\~ao \\
\midrule
Projeto (aloca\c{c}\~ao) & $1{,}08\!\times\!10^{-3}$ & $6{,}83\!\times\!10^{-4}$ & $5{,}4\!\times\!10^{-4}$ & overshoot $\sim$25\% (3.\textsuperscript{a} ordem) \\
``Melhorado'' (sim) & $1{,}08\!\times\!10^{-3}$ & $3{,}0\!\times\!10^{-4}$ & $9{,}0\!\times\!10^{-4}$ & menos $K_i$, mais $K_d$ ($\zeta\!\approx\!1{,}15$) \\
Default firmware & $8{,}0\!\times\!10^{-4}$ & $2{,}0\!\times\!10^{-5}$ & $1{,}2\!\times\!10^{-2}$ & conservador, manso no boot \\
\textbf{Em opera\c{c}\~ao} & $8{,}0\!\times\!10^{-4}$ & $2{,}0\!\times\!10^{-5}$ & $1{,}2\!\times\!10^{-2}$ & $=$ default do firmware; afinado no rig (com TRIM) \\
\bottomrule
\end{tabularx}
\end{table}

\begin{intuicao}[Por que tantos conjuntos diferentes?]
Os ganhos \emph{anal\'iticos} valem para o modelo ideal sem atraso. O hardware tem
atraso de sensor, ru\'ido e satura\c{c}\~ao, ent\~ao os ganhos de \emph{opera\c{c}\~ao}
s\~ao ajustados experimentalmente --- e funcionam apoiados em dois aliados do
firmware: a \textbf{rejei\c{c}\~ao de leituras esp\'urias} e o \textbf{TRIM}. Todos
saem da mesma teoria; mudam as prioridades (ideal vs.\ robusto).
\end{intuicao}

%==============================================================================
\chapter{$\zeta$ (Zeta) e $\omega_n$}\label{cap:zeta}
%==============================================================================
\index{zeta@$\zeta$}\index{omega n@$\omega_n$}

Todo sistema de 2.\textsuperscript{a} ordem padr\~ao escreve-se como
\begin{equation}
\frac{\omega_n^2}{s^2+2\zeta\omega_n s+\omega_n^2}.
\end{equation}
Dois n\'umeros resumem \emph{tudo} sobre a resposta: $\omega_n$ e $\zeta$.

\section{Frequ\^encia natural $\omega_n$}
\'E a velocidade ``de rel\'ogio'' do sistema: quanto maior $\omega_n$, mais
r\'apida a resposta. Na nossa malha PD, $\omega_n=\sqrt{A K_p}$.

\section{Coeficiente de amortecimento $\zeta$}
Mede \emph{quanto} a oscila\c{c}\~ao \'e atenuada:
\begin{center}
\begin{tabular}{ll}
$\zeta=0$ & oscila para sempre (n\~ao amortecido) \\
$0<\zeta<1$ & \textbf{subamortecido} (oscila e decai) \\
$\zeta=1$ & \textbf{criticamente amortecido} (mais r\'apido sem overshoot) \\
$\zeta>1$ & \textbf{sobreamortecido} (lento, sem overshoot) \\
\end{tabular}
\end{center}
Na malha PD, $\zeta=\dfrac{K_d}{2}\sqrt{\dfrac{A}{K_p}}$.

\begin{figure}[h]\centering
\begin{tikzpicture}
\begin{axis}[width=12.5cm,height=6.4cm,xlabel={$\omega_n t$ (adimensional)},
  ylabel={resposta ao degrau},xmin=0,xmax=14,ymin=0,ymax=1.8,
  legend pos=south east,grid=both,grid style={corcinza!20},
  legend style={font=\footnotesize},
  title={Fam\'ilia de respostas por $\zeta$},
  every axis title/.style={font=\small\bfseries,color=corprim}]
  \addplot[coralert,very thick,domain=0:14,samples=200]
    {1 - exp(-0.1*x)/sqrt(1-0.1^2)*sin(deg(sqrt(1-0.1^2)*x)+acos(0.1))};
  \addlegendentry{$\zeta=0{,}1$}
  \addplot[coracc,very thick,domain=0:14,samples=200]
    {1 - exp(-0.3*x)/sqrt(1-0.3^2)*sin(deg(sqrt(1-0.3^2)*x)+acos(0.3))};
  \addlegendentry{$\zeta=0{,}3$}
  \addplot[corroxo,very thick,domain=0:14,samples=200]
    {1 - exp(-0.7*x)/sqrt(1-0.7^2)*sin(deg(sqrt(1-0.7^2)*x)+acos(0.7))};
  \addlegendentry{$\zeta=0{,}7$}
  \addplot[corsec,very thick,domain=0:14,samples=200] {1-(1+x)*exp(-x)};
  \addlegendentry{$\zeta=1$}
  \draw[corcinza,dashed] (axis cs:0,1) -- (axis cs:14,1);
\end{axis}
\end{tikzpicture}
\caption{Quanto menor $\zeta$, maior o sobressinal e mais oscila\c{c}\~oes; $\zeta=1$ \'e o limiar sem overshoot.}
\label{fig:zeta}
\end{figure}

\begin{exemplo}[$\zeta$ dos conjuntos reais (parte PD)]
Com $A=7007$:
\begin{itemize}[leftmargin=1.4em]
  \item \textbf{Default} ($K_p=8\!\times\!10^{-4},K_d=1{,}2\!\times\!10^{-2}$):
  $\omega_n=\sqrt{7007\cdot8\!\times\!10^{-4}}=2{,}37$ rad/s,
  $\zeta=\tfrac{1{,}2\times10^{-2}}{2}\sqrt{7007/8\!\times\!10^{-4}}\approx17{,}8$ (muito lento).
  \item \textbf{Em opera\c{c}\~ao} ($K_p=6\!\times\!10^{-4},K_d=1\!\times\!10^{-4}$):
  $\omega_n=2{,}05$ rad/s, $\zeta\approx0{,}17$ (bem subamortecido).
\end{itemize}
Eis a leitura te\'orica do que voc\^e sente na mesa: com $\zeta\approx0{,}17$, a
resposta \'e r\'apida por\'em \emph{oscilat\'oria} --- exatamente o regime onde a
mesa ``surta'' sob qualquer perturba\c{c}\~ao se n\~ao houver as prote\c{c}\~oes
do firmware.
\end{exemplo}

\begin{projeto}[O que voc\^e \emph{v\^e} na mesa quando $\zeta$ e $\omega_n$ mudam]
\begin{itemize}[leftmargin=1.4em]
  \item \textbf{$\zeta$ baixo} (pouco $K_d$): a bola \textbf{passa do centro v\'arias
  vezes}, vai-e-volta, demora a assentar --- e ``surta'' a qualquer empurr\~ao.
  \item \textbf{$\zeta\approx1$} (bom $K_d$): a bola chega ao centro \textbf{firme,
  sem passar do ponto}. \'E o alvo da sintonia.
  \item \textbf{$\zeta$ alto} (muito $K_d$): a bola volta \textbf{devagar, arrastada},
  como se estivesse em mel.
  \item \textbf{$\omega_n$ maior} (mais $K_p$): \textbf{tudo fica mais r\'apido},
  corre\c{c}\~oes mais en\'ergicas. Mas como $\zeta=\tfrac{K_d}{2}\sqrt{A/K_p}$
  \emph{cai} quando $K_p$ sobe, aumentar $K_p$ \emph{sozinho} deixa a bola r\'apida
  \textbf{e} mais oscilat\'oria --- por isso $K_p$ e $K_d$ s\~ao ajustados juntos.
\end{itemize}
\end{projeto}

%==============================================================================
\chapter{Sobressinal (Overshoot)}\label{cap:overshoot}
%==============================================================================
\index{overshoot}\index{sobressinal}

\section{O que \'e}
\textbf{Overshoot} ($M_p$) \'e o quanto a resposta \emph{ultrapassa} o valor final
antes de assentar. Na mesa: a bola passa do centro e volta.

\section{Rela\c{c}\~ao com $\zeta$ e f\'ormula}
Para o sistema de 2.\textsuperscript{a} ordem subamortecido,
\begin{equation}
\boxed{\ M_p = \exp\!\left(\dfrac{-\zeta\pi}{\sqrt{1-\zeta^2}}\right)\times 100\%\ }
\label{eq:mp}
\end{equation}
Invertendo (\'util no projeto: do $PO$ desejado para o $\zeta$ necess\'ario):
\begin{equation}
\zeta = \frac{-\ln(M_p)}{\sqrt{\pi^2+\ln^2(M_p)}}.
\end{equation}

\begin{exemplo}[Verificando o requisito do projeto]
Com $\zeta=0{,}636$:
\[
M_p=\exp\!\left(\frac{-0{,}636\,\pi}{\sqrt{1-0{,}636^2}}\right)
=\exp\!\left(\frac{-1{,}998}{0{,}772}\right)=\exp(-2{,}589)\approx 0{,}075=7{,}5\%.
\]
Confere com o requisito $PO\le7{,}5\%$ usado para escolher $\zeta$ no Cap.~\ref{cap:ganhos}.
\end{exemplo}

\begin{table}[h]\centering
\caption{Overshoot em fun\c{c}\~ao de $\zeta$.}
\label{tab:mp}
\begin{tabular}{cc}
\toprule $\zeta$ & $M_p$ \\ \midrule
0,1 & 73\% \\ 0,3 & 37\% \\ 0,5 & 16\% \\ 0,636 & 7,5\% \\ 0,7 & 4,6\% \\ 1,0 & 0\% \\
\bottomrule
\end{tabular}
\end{table}

%==============================================================================
\chapter{Tempos de Resposta}\label{cap:tempos}
%==============================================================================
\index{tempo de acomodacao@tempo de acomoda\c{c}\~ao}

Tr\^es tempos caracterizam ``qu\~ao r\'apida'' \'e a resposta.

\begin{conceito}[Defini\c{c}\~oes]
\textbf{$t_r$ (subida):} tempo para ir de $\sim$10\% a $\sim$90\% do valor final.
\textbf{$t_p$ (pico):} instante do sobressinal m\'aximo.
\textbf{$t_s$ (acomoda\c{c}\~ao):} tempo para entrar (e ficar) numa faixa de
$\pm2\%$ em torno do valor final.
\end{conceito}

\begin{align}
t_p &= \frac{\pi}{\omega_d}=\frac{\pi}{\omega_n\sqrt{1-\zeta^2}}, &
t_r &\approx \frac{1{,}8}{\omega_n}, &
t_s &\approx \frac{4}{\zeta\omega_n}\ (\text{crit\'erio de }2\%).
\end{align}

\begin{figure}[h]\centering
\begin{tikzpicture}
\begin{axis}[width=12.5cm,height=6.4cm,xlabel={tempo (s)},ylabel={posi\c{c}\~ao},
  xmin=0,xmax=6,ymin=0,ymax=1.5,grid=both,grid style={corcinza!20},
  title={Marcos temporais ($\zeta=0{,}4$, $\omega_n=2{,}19$)},
  every axis title/.style={font=\small\bfseries,color=corprim}]
  \addplot[corprim,very thick,domain=0:6,samples=200]
    {1 - exp(-0.4*2.19*x)/sqrt(1-0.4^2)*sin(deg(2.19*sqrt(1-0.4^2)*x)+acos(0.4))};
  \draw[corsec,dashed] (axis cs:0,1.02) -- (axis cs:6,1.02);
  \draw[corsec,dashed] (axis cs:0,0.98) -- (axis cs:6,0.98);
  \draw[corcinza,dashed] (axis cs:0,1) -- (axis cs:6,1);
  \node[coracc,font=\footnotesize] at (axis cs:1.6,1.42) {$t_p$ (pico)};
  \draw[coracc,dotted,thick] (axis cs:1.56,0)--(axis cs:1.56,1.34);
  \node[corsec,font=\footnotesize,anchor=west] at (axis cs:4.1,1.1) {faixa $\pm2\%$};
\end{axis}
\end{tikzpicture}
\caption{Resposta ao degrau com $t_p$ e a faixa de acomoda\c{c}\~ao de $\pm2\%$.}
\label{fig:tempos}
\end{figure}

\begin{exemplo}[Tempos do projeto]
Com $\zeta=0{,}636$ e $\omega_n=2{,}19$: $\omega_d=2{,}19\sqrt{1-0{,}636^2}=1{,}69$ rad/s,
$t_p=\pi/1{,}69\approx1{,}86$ s, $t_s\approx4/(0{,}636\cdot2{,}19)\approx2{,}87$ s.
Coerente com o alvo $t_s\le2{,}5$--$3$ s.
\end{exemplo}

%==============================================================================
\chapter{M\'etodo de Ziegler--Nichols}\label{cap:zn}
%==============================================================================
\index{Ziegler-Nichols}

\section{A ideia}
Quando n\~ao se tem o modelo, Ziegler--Nichols (ZN) fornece ganhos a partir de
\textbf{um experimento}: leva-se o sistema ao limite da estabilidade.

\begin{conceito}[$K_u$ e $T_u$]
Com apenas o termo P ($K_i=K_d=0$), aumente $K_p$ at\'e a sa\'ida entrar em
\textbf{oscila\c{c}\~ao sustentada} (amplitude constante). Esse ganho \'e o
\textbf{ganho cr\'itico} $K_u$; o per\'iodo dessa oscila\c{c}\~ao \'e o
\textbf{per\'iodo cr\'itico} $T_u$.
\end{conceito}

\begin{table}[h]\centering
\caption{Tabela cl\'assica de Ziegler--Nichols (malha fechada).}
\label{tab:zn}
\begin{tabular}{lccc}
\toprule
Controlador & $K_p$ & $T_i$ & $T_d$ \\ \midrule
P   & $0{,}50\,K_u$ & --- & --- \\
PI  & $0{,}45\,K_u$ & $0{,}83\,T_u$ & --- \\
PID & $0{,}60\,K_u$ & $0{,}50\,T_u$ & $0{,}125\,T_u$ \\
\bottomrule
\end{tabular}
\end{table}
\noindent (Lembrando: $K_i=K_p/T_i$ e $K_d=K_p\,T_d$.)

\begin{figure}[h]\centering
\begin{tikzpicture}
\begin{axis}[width=12cm,height=4.6cm,xlabel={tempo (s)},ylabel={erro},
  xmin=0,xmax=9,ymin=-1.4,ymax=1.4,grid=both,grid style={corcinza!20},
  title={Oscila\c{c}\~ao sustentada no ganho cr\'itico $K_u$},
  every axis title/.style={font=\small\bfseries,color=corprim}]
  \addplot[coralert,very thick,domain=0:9,samples=220] {sin(deg(2*pi*x/2.9))};
  \draw[corsec,<->] (axis cs:2.0,1.15) -- node[above,font=\footnotesize]{$T_u$} (axis cs:4.9,1.15);
\end{axis}
\end{tikzpicture}
\caption{No limite, o erro oscila com amplitude constante e per\'iodo $T_u$.}
\label{fig:zn}
\end{figure}

\begin{projeto}[Como voc\^e faria o experimento de ZN na mesa]
\begin{enumerate}[leftmargin=1.5em]
  \item Pelo terminal: \texttt{PID <kp> 0 0} (zera I e D). Ligue \texttt{ERROR} ou
  \texttt{SHOW} para acompanhar.
  \item Suba $K_p$ aos poucos (\texttt{KP 0.0005}, \texttt{KP 0.0008}, \dots) at\'e a
  bola \textbf{oscilar com amplitude constante} em torno do centro. Esse \'e $K_u$.
  \item Me\c{c}a o \textbf{per\'iodo} dessa oscila\c{c}\~ao com cron\^ometro ou pelo
  \emph{timestamp} da telemetria: \'e $T_u$.
  \item Calcule $K_p=0{,}6K_u$, $T_i=0{,}5T_u$, $T_d=0{,}125T_u$ e converta para
  $K_i,K_d$. Use como \textbf{chute inicial} e refine (Cap.~\ref{cap:sintonia}).
\end{enumerate}
\end{projeto}

\begin{atencao}
ZN tende a dar respostas \textbf{agressivas} (overshoot $\sim$25\%). \'E excelente
ponto de partida, mas quase sempre conv\'em \textbf{aumentar $K_d$} (ou reduzir
$K_p$) depois para amortecer --- coerente com nossa planta sens\'ivel ao
derivativo.
\end{atencao}

%##############################################################################
\part{Controle Digital}
%##############################################################################

%==============================================================================
\chapter{Fundamentos de Controle Digital}\label{cap:digital}
%==============================================================================
\index{controle digital}\index{amostragem}

O ESP32 n\~ao resolve integrais cont\'inuas: ele executa o controlador em
\textbf{instantes discretos}, a cada $T$ segundos. Isso traz quatro conceitos.

\section{Amostragem e per\'iodo $T$}
A cada $T$ o firmware l\^e o sensor, calcula $u$ e comanda os motores.
\begin{projeto}
Nosso la\c{c}o roda a \textbf{100\,Hz} (\texttt{LOOP\_HZ = 100}), logo
$T=0{,}01$ s. Subimos de 50 para 100\,Hz justamente para \textbf{reduzir o atraso
de amostragem pela metade} e melhorar a resposta.
\end{projeto}

\section{Quantiza\c{c}\~ao}
\index{quantizacao@quantiza\c{c}\~ao}
O ADC da tela tem resolu\c{c}\~ao finita (12 bits, $0$--$4095$); a posi\c{c}\~ao em
mm vem em ``degraus''. Os motores tamb\'em s\~ao discretos (passos). Isso aparece
como pequenos ``micro-movimentos'' perto do centro --- limite f\'isico, n\~ao bug.

\section{Teorema de Nyquist e \emph{aliasing}}
\index{Nyquist}\index{aliasing}
\begin{conceito}[Nyquist]
Para reconstruir um sinal sem ambiguidade, a frequ\^encia de amostragem deve ser
\textbf{maior que o dobro} da maior frequ\^encia presente: $f_s > 2 f_{max}$. A
metade de $f_s$ \'e a \textbf{frequ\^encia de Nyquist}.
\end{conceito}
\textbf{Aliasing} \'e o que ocorre se essa regra \'e violada: componentes r\'apidas
``se disfar\c{c}am'' de lentas e contaminam a medida.

\begin{exemplo}[Folga de amostragem no nosso caso]
$f_s=100$ Hz $\Rightarrow$ Nyquist $=50$ Hz. A din\^amica da bola tem
$\omega_n\approx2{,}19$ rad/s, ou seja $f_n=\omega_n/2\pi\approx0{,}35$ Hz.
A amostragem \'e mais de \textbf{140 vezes} mais r\'apida que a din\^amica --- folga
enorme, e a aproxima\c{c}\~ao ``cont\'inua'' dos cap\'itulos anteriores vale com
seguran\c{c}a.
\end{exemplo}

%==============================================================================
\chapter{Atrasos no Sistema}\label{cap:atrasos}
%==============================================================================
\index{atraso}

Todo sistema real \textbf{demora} a reagir. Entre a bola se mover e a mesa
responder de fato, passam-se alguns milissegundos --- e atraso, em controle, \'e
\textbf{inimigo da estabilidade}. Vejamos de onde vem e por qu\^e.

\section{As tr\^es fontes de atraso}
\begin{enumerate}[leftmargin=1.5em]
  \item \textbf{Atraso de amostragem.} O firmware s\'o ``olha'' a bola a cada
  $T=10$ ms. Em m\'edia, a leitura usada j\'a tem \textbf{$T/2=5$ ms} de idade.
  \item \textbf{Atraso de processamento.} Ler o ADC, filtrar, rodar os dois PIDs e a
  cinem\'atica consome tempo de CPU; no nosso caso, pequeno perante os 10 ms do la\c{c}o.
  \item \textbf{Atraso do atuador.} O motor n\~ao salta para o \^angulo: ele
  \textbf{acelera e desacelera} (perfil trapezoidal). Entre comandar a inclina\c{c}\~ao
  e ela \emph{acontecer} passam-se alguns milissegundos.
\end{enumerate}

\begin{table}[h]\centering
\caption{Ordem de grandeza dos atrasos no Ball Balancer.}
\label{tab:atrasos}
\small
\begin{tabularx}{\textwidth}{l l X}
\toprule
Fonte & Ordem & Como o projeto mitiga \\
\midrule
Amostragem & $\sim$5 ms ($T/2$) & subir a taxa: 50 $\to$ \textbf{100 Hz} (metade do atraso) \\
Processamento & $<$1 ms & c\'odigo enxuto; telemetria limitada \\
Filtro (IIR/derivada) & $\sim$10--25 ms & $\alpha$ ajustado (ru\'ido $\times$ atraso) \\
Atuador (rampa) & poucos ms & acelera\c{c}\~ao alta o bastante, sem perder passo \\
\bottomrule
\end{tabularx}
\end{table}

\section{Por que atraso reduz a estabilidade}
\begin{intuicao}[Dirigir olhando pelo retrovisor]
Imagine corrigir a dire\c{c}\~ao de um carro vendo a estrada com 1 segundo de
\textbf{atraso}. Voc\^e reage ao que \emph{j\'a passou}: corrige tarde, passa do
ponto, corrige de novo --- e come\c{c}a a serpentear. O atraso transforma um
controle calmo em oscilat\'orio.
\end{intuicao}

Tecnicamente, o atraso \textbf{consome margem de fase}: o termo derivativo, que
deveria \emph{antecipar} o movimento, passa a agir sobre informa\c{c}\~ao
\emph{velha} e perde efic\'acia. Por isso n\~ao adianta s\'o ``encher de $K_d$'':
com muito atraso, nem o derivativo segura a oscila\c{c}\~ao.

\begin{projeto}[Como o projeto combate o atraso]
\begin{itemize}[leftmargin=1.4em]
  \item \textbf{La\c{c}o a 100 Hz} (em vez de 50): corta o atraso de amostragem pela
  metade --- melhoria direta de resposta.
  \item \textbf{Filtro derivativo com $\tau_d$ moderado}: corta ru\'ido sem
  adicionar atraso demais (Cap.~\ref{cap:filtro}).
  \item \textbf{Rejei\c{c}\~ao de \emph{outliers}} no lugar de filtragem pesada:
  remove o \emph{glitch} \textbf{sem} a lentid\~ao de um filtro forte.
\end{itemize}
\end{projeto}

%==============================================================================
\chapter{O Dom\'inio Z}\label{cap:z}
%==============================================================================
\index{transformada Z}\index{dominio z@dom\'inio $z$}

\section{O que \'e e por que existe}
Assim como Laplace ($s$) \'e a linguagem natural do tempo \emph{cont\'inuo}, a
\textbf{Transformada Z} ($z$) \'e a do tempo \emph{discreto}. Ela transforma
sequ\^encias $f[k]$ em fun\c{c}\~oes de $z$:
\begin{equation}
F(z)=\sum_{k=0}^{\infty} f[k]\,z^{-k}.
\end{equation}
O operador-chave \'e o \textbf{atraso de uma amostra}: $z^{-1}$ significa ``o valor
de uma amostra atr\'as''. Isso \'e tudo o que um microcontrolador precisa --
guardar valores anteriores.

\section{Rela\c{c}\~ao entre $s$ e $z$}
A ponte exata \'e
\begin{equation}
\boxed{\ z = e^{sT}\ }
\end{equation}
Ela mapeia o plano-$s$ no plano-$z$ assim:
\begin{center}
\begin{tabular}{ll}
semiplano esquerdo de $s$ (est\'avel) & $\to$ \textbf{interior} do c\'irculo unit\'ario de $z$ \\
eixo imagin\'ario de $s$ & $\to$ \textbf{c\'irculo} unit\'ario ($|z|=1$) \\
semiplano direito de $s$ (inst\'avel) & $\to$ \textbf{fora} do c\'irculo unit\'ario \\
\end{tabular}
\end{center}

\begin{figure}[h]\centering
\begin{tikzpicture}[scale=0.95]
  \begin{scope}
    \draw[seta] (-2.2,0)--(1.0,0) node[right,font=\footnotesize]{$\sigma$};
    \draw[seta] (0,-1.7)--(0,1.7) node[above,font=\footnotesize]{$j\omega$};
    \fill[corsec!12] (-2.2,-1.6) rectangle (0,1.6);
    \node[corsec!70!black,font=\footnotesize] at (-1.1,1.25) {est\'avel};
    \node[font=\small\bfseries,corprim] at (-0.6,-1.95) {plano-$s$};
  \end{scope}
  \begin{scope}[xshift=5.2cm]
    \draw[seta] (-1.8,0)--(1.8,0) node[right,font=\footnotesize]{$\mathrm{Re}$};
    \draw[seta] (0,-1.7)--(0,1.7) node[above,font=\footnotesize]{$\mathrm{Im}$};
    \fill[corsec!12] (0,0) circle (1.2);
    \draw[corprim,thick] (0,0) circle (1.2);
    \node[corsec!70!black,font=\footnotesize] at (0,0.6) {est\'avel};
    \node[font=\small\bfseries,corprim] at (0,-1.95) {plano-$z$ ($z=e^{sT}$)};
  \end{scope}
  \draw[seta,coracc,thick] (1.3,0.9) to[bend left=20] (3.0,0.9);
\end{tikzpicture}
\caption{O mapa $z=e^{sT}$ ``enrola'' o semiplano est\'avel de $s$ no interior do c\'irculo unit\'ario de $z$.}
\label{fig:zplane}
\end{figure}

\begin{intuicao}
Em controle cont\'inuo perguntamos ``os polos est\~ao \`a esquerda?''. Em controle
digital, a pergunta vira ``os polos est\~ao \textbf{dentro do c\'irculo
unit\'ario}?''. \'E o mesmo conceito de estabilidade, vestido para o tempo discreto.
\end{intuicao}

%==============================================================================
\chapter{Discretiza\c{c}\~ao}\label{cap:discret}
%==============================================================================
\index{discretizacao@discretiza\c{c}\~ao}

Para rodar o PID cont\'inuo no ESP32, aproximamos as opera\c{c}\~oes de
c\'alculo por f\'ormulas que usam apenas amostras. H\'a tr\^es m\'etodos cl\'assicos
de aproximar a integral/derivada (isto \'e, o operador $s$).

\begin{table}[h]\centering
\caption{M\'etodos de discretiza\c{c}\~ao (aproxima\c{c}\~oes de $s$).}
\label{tab:discret}
\small
\begin{tabularx}{\textwidth}{l l X}
\toprule
M\'etodo & Aproxima\c{c}\~ao de $s$ & Caracter\'istica \\
\midrule
Euler \emph{forward} (regressiva p/ frente) & $s\approx\dfrac{z-1}{T}$ &
Simples; pode instabilizar se $T$ grande. \\
Euler \emph{backward} (para tr\'as) & $s\approx\dfrac{z-1}{Tz}$ &
Sempre est\'avel; bom para a integral. \\
Tustin (bilinear/trapezoidal) & $s\approx\dfrac{2}{T}\dfrac{z-1}{z+1}$ &
Mais preciso; preserva melhor a fase. \\
\bottomrule
\end{tabularx}
\end{table}

\begin{conceito}[Vantagens e desvantagens]
\textbf{Forward} \'e a mais barata (s\'o usa o valor atual), mas a menos robusta.
\textbf{Backward} \'e incondicionalmente est\'avel --- por isso \'e segura para o
\emph{integrador}. \textbf{Tustin} \'e a mais fiel \`a din\^amica cont\'inua, ao
custo de um pouco mais de conta.
\end{conceito}

\begin{projeto}[O que o seu firmware usa]
O \texttt{pid.c} adota uma combina\c{c}\~ao pragm\'atica, comum em
microcontroladores:
\begin{itemize}[leftmargin=1.4em]
  \item \textbf{Integral por Euler \emph{forward}}: \texttt{p->integ += err*dt;}
  --- i.e.\ $I[k]=I[k-1]+e[k]\,T$.
  \item \textbf{Derivada por diferen\c{c}a regressiva}: $\dfrac{e[k]-e[k-1]}{T}$,
  seguida de um filtro passa-baixa (Cap.~\ref{cap:filtro}).
\end{itemize}
Com $T=0{,}01$ s e din\^amica lenta ($\omega_n\!\approx\!2{,}2$), qualquer um dos
m\'etodos daria praticamente o mesmo resultado --- a escolha \'e por simplicidade.
\end{projeto}

%==============================================================================
\chapter{O PID Discreto}\label{cap:pidd}
%==============================================================================
\index{PID discreto}

\section{Do cont\'inuo ao discreto, termo a termo}
Partindo de $u(t)=K_p e+K_i\!\int e+K_d\dot e$ e aplicando as aproxima\c{c}\~oes do
Cap.~\ref{cap:discret}:

\paragraph{Proporcional.} Direto: $P[k]=K_p\,e[k]$.

\paragraph{Integral (Euler forward).}
\begin{equation}
I[k]=I[k-1]+e[k]\,T, \qquad \text{termo} = K_i\,I[k].
\end{equation}

\paragraph{Derivativo (diferen\c{c}a regressiva).}
\begin{equation}
d[k]=\frac{e[k]-e[k-1]}{T}, \qquad \text{termo} = K_d\,d[k].
\end{equation}

\section{A equa\c{c}\~ao implementada}
Juntando, a lei de controle discreta que roda a cada 10\,ms \'e
\begin{equation}
\boxed{\ u[k]=K_p\,e[k]+K_i\,I[k]+K_d\,d_{lpf}[k]\ }
\label{eq:uk}
\end{equation}
seguida de \textbf{satura\c{c}\~ao} e \textbf{anti-windup} (Cap.~\ref{cap:windup}), onde
$d_{lpf}$ \'e a derivada \emph{filtrada} (Cap.~\ref{cap:filtro}). A cada borda de ``sem bola'', o
estado ($I$ e hist\'orico) \'e \textbf{zerado} (\texttt{pid\_reset}) para n\~ao
arrastar lixo de uma tentativa \`a pr\'oxima.

%==============================================================================
\chapter{Filtro Derivativo}\label{cap:filtro}
%==============================================================================
\index{filtro derivativo}\index{ruido@ru\'ido}

\section{O problema: o derivativo amplifica ru\'ido}
A derivada \emph{bruta} divide a diferen\c{c}a de amostras por $T=0{,}01$ s --- ou
seja, \textbf{multiplica por 100}. Um pequeno tremor de ru\'ido no sensor entre
duas amostras vira um pico enorme no termo D, que ``chuta'' os motores.

\section{A solu\c{c}\~ao: passa-baixa de 1.\textsuperscript{a} ordem}
Filtramos a derivada com um IIR (m\'edia exponencial):
\begin{equation}
d_{lpf}[k]=d_{lpf}[k-1]+\alpha\big(d[k]-d_{lpf}[k-1]\big),
\qquad \alpha=\frac{T}{\tau_d+T}.
\end{equation}
Isto equivale, no cont\'inuo, a $C_d(s)=\dfrac{K_d s}{1+\tau_d s}$: derivativo em
baixas frequ\^encias, ganho limitado em altas.

\begin{exemplo}[O $\alpha$ do firmware]
Com $\tau_d=0{,}03$ s e $T=0{,}01$ s:
\[
\alpha=\frac{0{,}01}{0{,}03+0{,}01}=0{,}25.
\]
O filtro atenua ru\'ido acima de $\sim\!\frac{1}{2\pi\tau_d}\approx5{,}3$ Hz,
preservando a din\^amica da bola (sub-Hz). A 50\,Hz, esse $\alpha$ era $0{,}4$;
ao dobrar a taxa para 100\,Hz, ele caiu para $0{,}25$ (mais suaviza\c{c}\~ao por
amostra, mesmo $\tau_d$).
\end{exemplo}

\begin{lstlisting}[caption={Derivada filtrada em \texttt{main/pid.c} (n\'ucleo).},label={lst:deriv}]
deriv = (err - p->prev_err) / dt;        /* diferenca regressiva    */
if (p->d_tau > 0.0f) {                   /* passa-baixa de 1a ordem */
    float a = dt / (p->d_tau + dt);      /* alpha = T/(tau_d+T)     */
    p->d_lpf += a * (deriv - p->d_lpf);
    deriv = p->d_lpf;
}
\end{lstlisting}

%==============================================================================
\chapter{Anti-Windup}\label{cap:windup}
%==============================================================================
\index{anti-windup}\index{saturacao@satura\c{c}\~ao}\index{windup}

\section{Satura\c{c}\~ao e \emph{windup}}
A inclina\c{c}\~ao \'e limitada a $\pm0{,}25$ rad. Se o erro \'e grande por muito
tempo, o \textbf{integrador acumula} um valor enorme (``\emph{windup}'') --- mas a
sa\'ida j\'a est\'a saturada e n\~ao melhora. Quando o erro finalmente inverte, o
integrador inflado demora a ``descarregar'', causando um \textbf{overshoot
gigante}.

\begin{figure}[h]\centering
\begin{tikzpicture}
\begin{axis}[width=12cm,height=4.8cm,xlabel={tempo},ylabel={sa\'ida $u$},
  xmin=0,xmax=6,ymin=0,ymax=0.55,grid=both,grid style={corcinza!20},
  legend style={font=\footnotesize},legend pos=north east,
  title={Com vs.\ sem anti-windup},
  every axis title/.style={font=\small\bfseries,color=corprim}]
  \addplot[coralert,very thick,domain=0:6,samples=120]
    {0.45*exp(-0.5*x)};
  \addlegendentry{sem anti-windup (estoura)}
  \addplot[corsec,very thick,domain=0:6,samples=120]
    {0.25*exp(-1.3*x)};
  \addlegendentry{com anti-windup}
  \draw[corcinza,dashed] (axis cs:0,0.25) -- (axis cs:6,0.25) node[right,font=\tiny]{$u_{max}$};
\end{axis}
\end{tikzpicture}
\caption{O anti-windup impede o integrador de ``inflar'' durante a satura\c{c}\~ao (ilustrativo).}
\label{fig:windup}
\end{figure}

\section{Back-calculation}
A t\'ecnica do firmware: quando a sa\'ida satura, \textbf{devolve-se ao integrador
exatamente o excesso}, de modo que ele nunca acumula al\'em do que produz a
satura\c{c}\~ao:
\begin{equation}
\text{se } u>u_{max}:\quad I \mathrel{-}= \frac{u-u_{max}}{K_i},\quad u=u_{max}.
\end{equation}

\begin{lstlisting}[caption={Anti-windup por \emph{back-calculation} em \texttt{main/pid.c}.},label={lst:aw}]
if (out > p->out_max) {
    if (p->ki != 0.0f) p->integ -= (out - p->out_max) / p->ki;
    out = p->out_max;
} else if (out < p->out_min) {
    if (p->ki != 0.0f) p->integ -= (out - p->out_min) / p->ki;
    out = p->out_min;
}
\end{lstlisting}

%##############################################################################
\part{Implementa\c{c}\~ao no Projeto}
%##############################################################################

%==============================================================================
\chapter{O Sensor: a Tela Resistiva}\label{cap:sensor}
%==============================================================================
\index{tela resistiva}\index{sensor}

Sem medir a bola, n\~ao h\'a realimenta\c{c}\~ao. O sensor do projeto \'e uma
\textbf{tela resistiva de 4 fios} ($187\times141$ mm) sobre a qual a bola repousa.

\section{Como funciona}
A tela tem \textbf{duas camadas} condutoras separadas por micro-espa\c{c}adores. O
peso da bola \textbf{encosta} uma camada na outra no ponto de contato. Para ler o
eixo X, aplica-se uma tens\~ao ao longo de uma camada (um \textbf{divisor de
tens\~ao} na dire\c{c}\~ao X) e l\^e-se a tens\~ao no ponto de contato pela outra
camada --- ela \'e \textbf{proporcional \`a posi\c{c}\~ao}. Trocando o papel das
camadas, l\^e-se Y. Dois ciclos $\Rightarrow$ coordenadas $(x,y)$.

\begin{intuicao}
A tela \'e um \textbf{potenci\^ometro 2D}: a posi\c{c}\~ao da bola vira uma
fra\c{c}\~ao da tens\~ao aplicada. Barato, simples e sem c\^amera.
\end{intuicao}

\section{Convers\~ao para coordenadas}
O ADC do ESP32 l\^e a tens\~ao como um n\'umero ``cru'' de 12 bits ($0$--$4095$). A
calibra\c{c}\~ao mapeia o cru para mil\'imetros por uma reta:
\[
x_{mm} = \frac{\text{cru}_x - \text{cru}_{\min}}{\text{cru}_{\max}-\text{cru}_{\min}}\times 187.
\]

\section{Calibra\c{c}\~ao}
O firmware aprende $\text{cru}_{\min}$ e $\text{cru}_{\max}$ de cada eixo colocando a
bola nos \textbf{4 cantos} (comando \texttt{CAL}); o \textbf{centro \'e derivado}
deles, n\~ao medido (medi-lo ``no olho'' enviesaria a reta). Os valores ficam no
NVS e valem ap\'os reiniciar.

\section{Ru\'ido e quantiza\c{c}\~ao}
\begin{itemize}[leftmargin=1.4em]
  \item \textbf{Ru\'ido el\'etrico}: motores e o acoplamento capacitivo entre as
  camadas poluem a leitura. Defesas: m\'edia de 8 amostras de ADC, filtro IIR
  ($\alpha=0{,}45$), \textbf{detec\c{c}\~ao digital} de presen\c{c}a (imune ao
  acoplamento) e \textbf{rejei\c{c}\~ao de \emph{outliers}}.
  \item \textbf{Quantiza\c{c}\~ao}: 12 bits em 187 mm d\~ao $\sim$0{,}05 mm por
  degrau de ADC --- fino o bastante; o limite pr\'atico \'e o ru\'ido, n\~ao a
  resolu\c{c}\~ao.
\end{itemize}

\section{Impacto no controle}
\begin{atencao}[O sensor \'e o teto do desempenho]
Ru\'ido e atraso de leitura \textbf{limitam o quanto de $K_d$} se pode usar: como o
derivativo amplifica ru\'ido (Cap.~\ref{cap:filtro}), uma leitura suja obriga a
filtrar mais --- o que \emph{adiciona} atraso (Cap.~\ref{cap:atrasos}). Por isso
investir na \textbf{qualidade da medida} (detec\c{c}\~ao digital, rejei\c{c}\~ao de
\emph{glitch}) foi t\~ao decisivo quanto sintonizar o PID: melhor sensor
$\Rightarrow$ mais $K_d$ \'util $\Rightarrow$ resposta mais firme.
\end{atencao}

%==============================================================================
\chapter{An\'alise do C\'odigo}\label{cap:codigo}
%==============================================================================
\index{firmware}

Aqui a teoria encontra o C. Tr\^es arquivos formam o cora\c{c}\~ao do sistema.

\section{\texttt{pid.c} --- o controlador por eixo}
A fun\c{c}\~ao \texttt{pid\_update} \'e a Eq.~\eqref{eq:uk} em c\'odigo, na ordem:
derivada filtrada $\to$ integral $\to$ soma $\to$ satura\c{c}\~ao/anti-windup.
\begin{lstlisting}[caption={N\'ucleo de \texttt{pid\_update} (resumido).},label={lst:pid}]
float err = meas - setpoint;             /* e[k] = medida - alvo    */
/* derivada filtrada (Cap. 18) */
float deriv = 0.0f;
if (p->primed) {
    deriv = (err - p->prev_err) / dt;
    float a = dt / (p->d_tau + dt);
    p->d_lpf += a * (deriv - p->d_lpf);
    deriv = p->d_lpf;
}
p->prev_err = err;
p->integ += err * dt;                    /* integral Euler forward  */
float out = p->kp*err + p->ki*p->integ + p->kd*deriv;  /* u[k] (Eq.) */
/* saturacao + anti-windup (Cap. 19) ... */
\end{lstlisting}
\begin{conceito}[Mapeamento teoria $\to$ c\'odigo]
\texttt{p->kp*err} $=K_p e[k]$;\quad
\texttt{p->ki*p->integ} $=K_i I[k]$;\quad
\texttt{p->kd*deriv} $=K_d d_{lpf}[k]$.\quad
A estrutura \texttt{pid\_t} guarda $K_p,K_i,K_d$, os limites de sa\'ida, o
acumulador \texttt{integ}, o erro anterior \texttt{prev\_err} e o estado do filtro
\texttt{d\_lpf}.
\end{conceito}

\section{\texttt{control.c} --- o la\c{c}o de 100\,Hz}
A cada 10\,ms: l\^e o toque (com detec\c{c}\~ao digital e rejei\c{c}\~ao de
\emph{outliers}), aplica o \textbf{debounce} de presen\c{c}a, roda os dois PIDs,
soma o \textbf{TRIM} e converte a inclina\c{c}\~ao em \^angulos.
\begin{lstlisting}[caption={Trecho central de \texttt{control.c} (sa\'ida do PID $\to$ motores).},label={lst:control}]
float ox = pid_update(&s_pid_x, p.x_mm, SETPOINT_X_MM, dt_real);
float oy = pid_update(&s_pid_y, p.y_mm, SETPOINT_Y_MM, dt_real);
nx = s_sign_x * ox;     /* realimentacao negativa: TILT_SIGN = -1  */
ny = s_sign_y * oy;
nx += s_trim_x;         /* feedforward de nivel (base torta)        */
ny += s_trim_y;
apply_tilt(nx, ny, thoff);   /* cinematica inversa 3RPS -> 3 motores */
\end{lstlisting}

\section{\texttt{kinematics.c} --- a cinem\'atica inversa 3RPS}
A fun\c{c}\~ao \texttt{machine\_theta} recebe o vetor normal $(n_x,n_y,1)$ da mesa e
devolve o \textbf{\^angulo de cada motor}. Ela normaliza a normal, posiciona a junta
esf\'erica de cada perna e resolve o tri\^angulo bra\c{c}o--biela por lei dos
cossenos (da\'i os dois \texttt{acos}):
\begin{lstlisting}[caption={Esqueleto de \texttt{machine\_theta} (perna A).},label={lst:kin}]
double nmag = sqrt(nx*nx + ny*ny + 1.0); /* normaliza (nx,ny,1)     */
nx /= nmag; ny /= nmag; double nz = 1.0/nmag;
/* ... posiciona a junta da perna (y, z) a partir da normal ... */
mag = sqrt(y*y + z*z);
angle = acos(y/mag)                                   /* geometria  */
      + acos((mag*mag + f*f - g*g)/(2.0*mag*f));      /* lei dos cossenos */
return angle * (180.0/M_PI);              /* radianos -> graus       */
\end{lstlisting}
\begin{projeto}[Do \^angulo para passos]
A sa\'ida em graus vira passos pelo fator
$\dfrac{200\cdot\text{microsteps}}{360^\circ}$. Com microstepping $1/16$:
$\dfrac{200\cdot16}{360}\approx8{,}9$ passos por grau. \'E o elo final entre o
``mundo cont\'inuo'' do controle e os passos discretos do motor.
\end{projeto}

%==============================================================================
\chapter{Arquitetura Completa do Sistema}\label{cap:arq}
%==============================================================================
\index{arquitetura}

\section{O fluxo de ponta a ponta}
\begin{figure}[h]\centering
\begin{tikzpicture}[node distance=8mm and 9mm,every node/.style={font=\footnotesize}]
  \node[blocoS] (sensor) {Tela resistiva\\(detec\c{c}\~ao digital\\+ rejei\c{c}\~ao de outliers)};
  \node[bloco,right=of sensor] (pid) {2$\times$ PID\\$u[k]$ (Eq.~\eqref{eq:uk})};
  \node[bloco,right=of pid] (trim) {$+$ TRIM\\(feedforward)};
  \node[bloco,right=of trim] (kin) {Cinem\'atica\\3RPS};
  \node[bloco,below=of kin] (step) {Steppers\\(perfil trapez.)};
  \node[blocoS,left=of step] (plat) {Plataforma\\(inclina\c{c}\~ao $\alpha$)};
  \node[blocoS,left=of plat] (ball) {Bola\\$A/s^2$};
  \draw[seta] (sensor) -- node[above]{$x,y$} (pid);
  \draw[seta] (pid) -- (trim);
  \draw[seta] (trim) -- node[above]{$n_x,n_y$} (kin);
  \draw[seta] (kin) -- node[right]{$\theta_{A,B,C}$} (step);
  \draw[seta] (step) -- (plat);
  \draw[seta] (plat) -- (ball);
  \draw[seta] (ball) -| node[pos=0.25,below]{realimenta\c{c}\~ao} (sensor);
\end{tikzpicture}
\caption{Arquitetura completa: o anel fecha do sensor \`a bola e volta, a 100\,Hz.}
\label{fig:arq}
\end{figure}

\section{Camadas de software}
\begin{table}[h]\centering
\caption{M\'odulos do firmware e seu papel.}
\label{tab:modulos}
\small
\begin{tabularx}{\textwidth}{>{\ttfamily\bfseries}l X}
\toprule
main.c & Inicializa\c{c}\~ao; seleciona modo CONTROL/MAPPING. \\
control.c & La\c{c}o de 100\,Hz, \emph{homing}, parser serial, TRIM, telemetria. \\
pid.c & PID por eixo (derivada filtrada + anti-windup). \\
kinematics.c & Cinem\'atica inversa 3RPS (\texttt{machine\_theta}). \\
steppers.c & Gera\c{c}\~ao de passos STEP/DIR (perfil trapezoidal, 10\,kHz). \\
touch\_screen.c & Leitura resistiva, detec\c{c}\~ao digital, filtro, outliers. \\
cal\_store.c & Persist\^encia em NVS (PID, calibra\c{c}\~ao, curso, ZERO, TRIM). \\
\bottomrule
\end{tabularx}
\end{table}

\section{O atuador: perfil trapezoidal}
\index{perfil trapezoidal}
Cada motor segue a posi\c{c}\~ao-alvo com um \textbf{perfil trapezoidal de
velocidade}: acelera at\'e \texttt{SPEED\_BALANCE}, navega, e desacelera a tempo de
parar \emph{exatamente} no alvo (velocidade segura $v_{lim}=\sqrt{2ad}$). Sem essa
rampa, a partida brusca causaria \textbf{perda de passos}. A banda do atuador
(centenas de Hz) \'e muito maior que a da bola ($\sim$0,35 Hz), por isso ele
\textbf{n\~ao compromete} a estabilidade projetada e pode ser tratado como ganho
quase-est\'atico na an\'alise (Cap.~\ref{cap:blocos}).

%==============================================================================
\chapter{Guia Pr\'atico de Sintonia}\label{cap:sintonia}
%==============================================================================
\index{sintonia}\index{tuning}

Todo o ajuste \'e feito \textbf{ao vivo pelo terminal serial} (USB-Serial-JTAG),
sem recompilar. Ligue \texttt{ERROR} para ver $e_x,e_y$ e a dist\^ancia ao centro a
cada segundo.

\section{Roteiro recomendado}
\begin{enumerate}[leftmargin=1.5em]
  \item \textbf{Ajuste $K_p$:} com $K_i=K_d=0$, suba $K_p$ at\'e a bola reagir bem,
  por\'em \emph{oscilando} um pouco. (\texttt{KP 0.0008}\dots)
  \item \textbf{Ajuste $K_d$:} adicione derivativo at\'e a oscila\c{c}\~ao
  \textbf{sumir} (alvo: $\zeta\!\approx\!1$, sem overshoot). \emph{\'E o passo que
  estabiliza} o duplo integrador. (\texttt{PID - - 0.012})
  \item \textbf{Ajuste $K_i$:} um \textbf{toque} para eliminar o desvio residual da
  mesa torta. (\texttt{PID - 0.0001 -})
  \item \textbf{TRIM:} com a bola centrada, \texttt{TRIM} e \texttt{TRIM SAVE}
  (Cap.~\ref{cap:futuro}).
  \item \textbf{Grave:} \texttt{PID SAVE}.
\end{enumerate}

\section{Reconhecendo cada patologia}
\begin{table}[h]\centering
\caption{Diagn\'ostico r\'apido por sintoma.}
\label{tab:diag}
\small
\begin{tabularx}{\textwidth}{l X l}
\toprule
Sintoma & Causa prov\'avel & A\c{c}\~ao \\
\midrule
\textbf{Oscila\c{c}\~ao} crescente & $K_p$ alto / $K_d$ baixo & subir $K_d$ ou baixar $K_p$ \\
\textbf{Instabilidade} (``surta'') & $\zeta$ muito baixo & subir $K_d$; checar outliers \\
\textbf{Lentid\~ao} & $K_p$ baixo / $K_d$ alt\'issimo & subir $K_p$ / baixar $K_d$ \\
\textbf{Desvio} fixo do centro & $K_i$ fraco / base torta & subir $K_i$; usar \texttt{TRIM} \\
\textbf{Satura\c{c}\~ao} constante & ganhos altos demais & reduzir; conferir anti-windup \\
\textbf{Trem\^or}/micro-movimento & ru\'ido / $K_d$ alto & filtrar mais; baixar $K_d$ \\
\bottomrule
\end{tabularx}
\end{table}

\subsection{Assinaturas visuais (o que olhar na bancada)}
\begin{itemize}[leftmargin=1.4em]
  \item \textbf{Excesso de $K_p$:} a bola fica \textbf{nervosa} --- vibra/oscila em
  torno do centro com amplitude que n\~ao morre. \emph{Baixe $K_p$ ou suba $K_d$.}
  \item \textbf{Falta de $K_d$:} a bola \textbf{passa do centro} e volta v\'arias
  vezes (overshoot grande) antes de assentar; a qualquer empurr\~ao, ``surta''.
  \emph{Suba $K_d$.}
  \item \textbf{Excesso de $K_i$:} \textbf{oscila\c{c}\~ao lenta e larga} (vai longe
  de um lado, depois do outro), com per\'iodo de v\'arios segundos --- \'e o
  \emph{windup}. \emph{Baixe $K_i$.}
  \item \textbf{Falta de $K_i$:} a bola estabiliza, por\'em \textbf{deslocada} do
  centro (mesa torta). \emph{Suba $K_i$ ou use \texttt{TRIM}.}
\end{itemize}

\begin{atencao}[A li\c{c}\~ao do seu pr\'oprio rig]
Quando o $K_d$ salvo ficou cerca de 120$\times$ abaixo do necess\'ario, a mesa
``surtava do nada'' --- $\zeta\!\approx\!0{,}17$ (Cap.~\ref{cap:zeta}). A cura foi exatamente a
teoria: \textbf{mais amortecimento} (subir $K_d$) e \textbf{rejei\c{c}\~ao de
leituras esp\'urias}. Confie no $\zeta$: ele prev\^e o comportamento antes de voc\^e
tocar na mesa.
\end{atencao}

\section{Teoria vs.\ Hardware: por que os ganhos n\~ao batem}
\index{teoria vs hardware}
Os ganhos calculados por aloca\c{c}\~ao de polos (Cap.~\ref{cap:ganhos}) quase nunca
s\~ao os ganhos \emph{finais} no rig --- e isso \textbf{n\~ao \'e um erro}. O modelo
$P(s)=A/s^2$ \'e ideal; o hardware carrega imperfei\c{c}\~oes que o modelo ignora.

\begin{table}[h]\centering
\caption{Por que o ganho de bancada difere do ganho de projeto.}
\label{tab:tvh}
\small
\begin{tabularx}{\textwidth}{l X}
\toprule
Fator real & Efeito sobre os ganhos \\
\midrule
\textbf{Ru\'ido} do sensor & limita $K_d$ (o derivativo amplifica ru\'ido) \\
\textbf{Satura\c{c}\~ao} ($\pm0{,}25$ rad) & ganhos altos saturam mais cedo; \emph{n\~ao} aceleram \\
\textbf{Atrasos} (amostragem, motor) & consomem margem de fase; pedem cautela com $K_p$ \\
\textbf{Limites mec\^anicos} (folga, atrito, base torta) & criam erros que o modelo ideal n\~ao prev\^e \\
\textbf{Cinem\'atica n\~ao linear} & o ganho efetivo varia com a inclina\c{c}\~ao \\
\bottomrule
\end{tabularx}
\end{table}

\begin{intuicao}[A teoria acerta o \emph{ponto de partida}, n\~ao o ponto final]
A aloca\c{c}\~ao de polos diz \emph{onde come\c{c}ar} e \emph{para que lado andar}:
se oscila, falta $K_d$; se desvia, falta $K_i$. O valor exato sai do ajuste fino no
hardware. \'E como uma receita: as propor\c{c}\~oes v\^em do livro, o ``ponto'' do
sal voc\^e acerta provando.
\end{intuicao}

\begin{projeto}[Preparado para a banca]
Se perguntarem \emph{``por que seus ganhos n\~ao s\~ao os calculados?''}, a resposta
\'e: o modelo \'e \textbf{linear, ideal e sem atraso}; o rig tem ru\'ido,
satura\c{c}\~ao, atraso e folga mec\^anica. A teoria deu o \textbf{projeto inicial}
(estrutura do controlador, ordem de grandeza e dire\c{c}\~ao de ajuste) e validou o
comportamento em simula\c{c}\~ao; o conjunto final foi \textbf{refinado
experimentalmente}, apoiado na rejei\c{c}\~ao de ru\'ido e no \emph{feedforward}
TRIM. Teoria e pr\'atica n\~ao se contradizem --- complementam-se.
\end{projeto}

%==============================================================================
\chapter{Melhorias Futuras}\label{cap:futuro}
%==============================================================================
\index{feedforward}\index{LQR}\index{MPC}\index{Kalman}

O PID resolve muito bem o Ball Balancer, mas a teoria moderna oferece extens\~oes.

\section{Feedforward (j\'a presente: o TRIM)}
\index{TRIM}
A base torta \'e uma \textbf{perturba\c{c}\~ao constante}. Em vez de esperar o
integral lento descobri-la a cada bola, o \textbf{TRIM} mede uma vez a inclina\c{c}\~ao
que o integral encontrou e a aplica como \emph{feedforward} fixo (salvo no NVS).
\begin{equation}
n_x \leftarrow \underbrace{s_x\,u_x}_{\text{PID}} + \underbrace{\textsc{trim}_x}_{\text{feedforward}}
\end{equation}
\begin{intuicao}
Feedforward \'e ``saber a resposta de antem\~ao''. Como a mesa est\'a sempre torta
do mesmo jeito, n\~ao faz sentido redescobrir isso a cada bola: aplicamos a
corre\c{c}\~ao direto, e o PID s\'o cuida do que \emph{varia}. \'E separar a
rejei\c{c}\~ao da perturba\c{c}\~ao DC da regula\c{c}\~ao em torno dela.
\end{intuicao}

\section{Para al\'em do PID (men\c{c}\~ao breve)}
O PID resolve muito bem este sistema. Se um dia se quiser ir adiante, ficam as
\textbf{cita\c{c}\~oes} a seguir --- \emph{fora do escopo} desta apostila, que foca
no dom\'inio completo do PID digital atual:
\begin{itemize}[leftmargin=1.4em]
  \item \textbf{Espa\c{c}o de estados / LQR:} projeta o controlador tratando
  posi\c{c}\~ao e velocidade como ``estado'', de forma \'otima e sistem\'atica ---
  generaliza a aloca\c{c}\~ao de polos do Cap.~\ref{cap:ganhos}.
  \item \textbf{Filtro de Kalman:} estima a velocidade de forma \'otima a partir da
  medida ruidosa, em vez da derivada num\'erica (Cap.~\ref{cap:filtro}).
  \item \textbf{MPC:} otimiza a a\c{c}\~ao respeitando \emph{explicitamente} a
  satura\c{c}\~ao de $\pm0{,}25$ rad.
  \item \textbf{Controle adaptativo:} reajusta os ganhos sozinho conforme a planta
  muda --- a fam\'ilia do \texttt{PIDAUTO} embarcado.
\end{itemize}
Todas \textbf{complementam}, n\~ao substituem, o PID consolidado aqui.

\begin{table}[h]\centering
\caption{Quando valeria a pena cada t\'ecnica avan\c{c}ada.}
\label{tab:avancado}
\small
\begin{tabularx}{\textwidth}{l X}
\toprule
T\'ecnica & Ganho potencial sobre o PID atual \\
\midrule
Feedforward (TRIM) & \textbf{j\'a implementado}: centra na hora, sem rampa do integral. \\
LQR / espa\c{c}o de estados & Sintonia sistem\'atica e \'otima; \'util se acoplar X--Y. \\
Kalman & Velocidade limpa $\Rightarrow$ pode subir $K_d$ sem ru\'ido. \\
MPC & Trata satura\c{c}\~ao e restri\c{c}\~oes de forma \'otima. \\
Adaptativo & Robustez a mudan\c{c}as (peso da bola, atrito). \\
\bottomrule
\end{tabularx}
\end{table}

\begin{projeto}[Mensagem final]
Voc\^e tem, hoje, um sistema \textbf{completo e funcional}: planta modelada,
ganhos com base te\'orica, PID digital com filtro e anti-windup, cinem\'atica 3RPS,
rejei\c{c}\~ao de ru\'ido e \emph{feedforward} para a base torta. As ``melhorias
futuras'' n\~ao s\~ao conserto --- s\~ao os \emph{pr\'oximos cap\'itulos} naturais
de quem j\'a domina o b\'asico. Boa defesa de projeto!
\end{projeto}

%==============================================================================
%  REFERENCIAS / INDICE
%==============================================================================
\backmatter
\nocite{*}
\bibliographystyle{ieeetr}
\bibliography{referencias}
\addcontentsline{toc}{chapter}{Refer\^encias}

\printindex

\end{document}
